% LaTeX master Datei(en) zusammengestellt von Falk-Jonatan Strube zur Nutzung an der Hochschule für Technik und Wirtschaft Dresden: https://github.com/genericFJS/htw
\documentclass{scrartcl}
\gdef\pathtomaster{../_LaTeX_master}
%% HTW Dresden Corporate Design
%% Copyright (C) Falk-Jonatan Strube <falkjonatanstrube@gmail.com>, 2017
%%% === Anfang Paket htwscr ===
\makeatletter
\ifcsname KOMAoption\endcsname%
\KOMAoption{footheight}{11mm}	% mehr Platz für Wasserzeichen o.ä. in der Fußzeile
\recalctypearea	% s.o.
\RequirePackage{geometry}
\geometry{top=25mm, bottom=25mm, headsep=10mm, inner=25mm, outer=25mm, footskip=15mm, marginparsep=2mm}
\RequirePackage{helvet}	% Corporate Design schreibt Arial/Helvetica vor
\renewcommand{\familydefault}{\sfdefault}	% s.o.
\RequirePackage{graphicx}	% für Grafiken
\RequirePackage{xcolor}		% für Farben
\definecolor{htworange}{RGB}{249,155,28}	% Corporate Design orange CMYK(0/45/100/0)
\definecolor{htwblue}{RGB}{0,116,188}			% Corporate Design blau CMYK(90/50/0/0)
\definecolor{htwgrey}{RGB}{128,128,128}		% Corporate Design gray 60% Schwarz [Annahme: CMYK(0,0,0,60)]
\RequirePackage{ifthen,xifthen}	% für Bedingungen/Optionen
\RequirePackage{etoolbox}				% s.o.

%%% Format des HTW-Headers (wird hier als Texteelement verwenden) nach dem Style der PowerPoint-Vorlagen. Weiterhin inspiriert durch das Paket tudscr (Corporate Design der TU Dresden)
%%% Dateipfad des Logos abhängig von der Verwendung als Paket/Input
\ifdefined\@isHTWCDpackage%
	\gdef\@pathtoimage{HTW}%
\else%
	\ifdefined\pathtomaster%
		\gdef\@pathtoimage{\pathtomaster /HTW}%
	\else 
		\gdef\@pathtoimage{HTW}%
	\fi%
\fi%
\newcommand{\pathtologo}[1]{\gdef\@pathtoimage{#1}}
%%% Berechnung der Größe des Logos
\newlength{\htwLogoHeight}
\setlength{\htwLogoHeight}{12mm}	% Mindesthöhe des Logos laut Corporate Design: 11mm
\newlength{\@htwLogoWidth}
\AtBeginDocument{
\def\printLogo{\includegraphics[height=\htwLogoHeight]{\@pathtoimage}}
\setlength\@htwLogoWidth{\widthof{\printLogo}}
}
%%% Definition des Header
\newcommand{\@htwheader}{{\noindent%
%%% Header ist abhängig von dem Vorhandensein von department, institute und chair
\newcommand{\hasDepInsCh}{\gdef\@hasDepInsCh{}}%
\ifdefined\@department\hasDepInsCh\fi%
\ifdefined\@institute\hasDepInsCh\fi%
\ifdefined\@chair\hasDepInsCh\fi%
\begin{minipage}[b][\htwLogoHeight]{\dimexpr \textwidth - \@htwLogoWidth - 2mm \relax}% 2mm Abstand zwischen Logo und Linie
{\small{\bfseries%
\ifdefined\@faculty\@faculty{}\fi} %
\strut}%
{\color{htworange}\hrule height 1pt}%
\vspace{.1mm}%
\end{minipage}
\hfill%
\begin{minipage}[b][\htwLogoHeight]{\@htwLogoWidth}
\includegraphics[height=\htwLogoHeight]{\@pathtoimage}
\end{minipage}
\ifdefined\@hasDepInsCh%
\\%
\begin{minipage}{\textwidth}
\small%
\ifdefined\@department\@department{}\fi%
\ifdefined\@institute\ifdefined\@department,~\fi \@institute\fi%
\ifdefined\@chair\ifdefined\@department,~\else\ifdefined\@institute,~\fi\fi\@chair\fi\strut
%{\color{htworange}\hrule height 0.5pt}
\end{minipage}
\fi
}}

%%% Definition der Fußzeilen
\usepackage{scrlayer-scrpage}
\@ifclassloaded{scrreprt}{%
\renewcommand*{\partpagestyle}{scrheadings}
\renewcommand*{\chapterpagestyle}{scrheadings}%
}{}
\pagestyle{scrheadings}
\cfoot{\vfill\pagemark}
\chead{}

%%% Einstellungen
\newcommand{\simpledate}{\gdef\@simpledate{}}		% Datum nicht mit "Eingereicht am:" Vorsatz
\newcommand{\nouppercase}{\gdef\@nouppercase{}}	% Titel/Parts/Chapters usw. nicht in Großbuchstaben

%%% Einiges erst zu Beginn des Dokumentes ausführen, um Einstellungen zu berücksichtigen
\AtBeginDocument{
\renewcommand\SS{SS}	% fix big ß
%%% Sektionen usw. ggf. in Großbuchstaben:
\providecommand{\chapterlinesformat}[3]{}	% compatibility
\providecommand{\chapterlineswithprefixformat}[3]{}	% compatibility
\providecommand{\sectionlinesformat}[4]{}	% compatibility
\ifdefined\@nouppercase\relax\else
\renewcommand\sectionlinesformat[4]{\@hangfrom{\hskip#2 #3}{\MakeUppercase{#4}}}
\@ifclassloaded{scrreprt}{%
\renewcommand\chapterlinesformat[3]{\@hangfrom{#2}{\MakeUppercase{#3}}}
\renewcommand\chapterlineswithprefixformat[3]{\MakeUppercase{#2#3}}%
}{}
\fi
%%% Anpassung der Part-Seite
\renewcommand*{\partheadstartvskip}{\@htwheader\vspace*{30mm}}
\renewcommand*{\raggedpart}{\raggedright}
\renewcommand*{\partformat}{\huge\ifdefined\@nouppercase \partname~\thepart\else\MakeUppercase{\partname~\thepart}\fi.}
\ifdefined\@nouppercase\setkomafont{part}{\Huge}\else\setkomafont{part}{\Huge\MakeUppercase}\fi
}

\newcommand{\bilingual}[3]{\ifdefined #1 \let #1\undefined\fi\newcommand{#1}{\iflanguage{english}{#3}{#2}}}	% Übersetzung von einigen Variablen in mehrere Sprachen

\newcommand{\checkThesis}[2]{\ifthenelse{\equal{#1}{diss}\OR\equal{#1}{doctoral}\OR\equal{#1}{phd}}{
\gdef#2{\dissertationname}}{
\ifthenelse{\equal{#1}{diploma}}{
\gdef#2{\diplomathesisname}}{
\ifthenelse{\equal{#1}{master}}{
\gdef#2{\masterthesisname}}{
\ifthenelse{\equal{#1}{bachelor}}{
\gdef#2{\bachelorthesisname}}{
\ifthenelse{\equal{#1}{student}}{
\gdef#2{\studentthesisname}}{
\ifthenelse{\equal{#1}{evidence}}{
\gdef#2{\studentresearchname}}{
\ifthenelse{\equal{#1}{project}}{
\gdef#2{\projectpapername}}{
\ifthenelse{\equal{#1}{seminar}}{
\gdef#2{\seminarpapername}}{
\ifthenelse{\equal{#1}{term}}{
\gdef#2{\termpapername}}{
\ifthenelse{\equal{#1}{research}}{
\gdef#2{\researchname}}{
\ifthenelse{\equal{#1}{log}}{
\gdef#2{\logname}}{
\ifthenelse{\equal{#1}{report}}{
\gdef#2{\reportname}}{
\ifthenelse{\equal{#1}{internship}}{
\gdef#2{\internshipname}}{
\ifthenelse{\equal{#1}{lecture}\OR\equal{#1}{lesson}\OR\equal{#1}{pract}}{
\bilingual{\professorname}{Vorlesung von}{Lecture by}
\simpledate
\gdef\@islecture{\relax}
\ifthenelse{\equal{#1}{lecture}}{\gdef#2{\lecturename}}{
\ifthenelse{\equal{#1}{lesson}}{\gdef#2{\lessonname}}{
\gdef#2{\practname}
}}}{
\gdef#2{#1}}}}}}}}}}}}}}}}

%%% Textfragmente der Titelseite in deutsch und englisch
\bilingual{\dateofbirthname}{Geboren am:}{Born on:}
\bilingual{\placeofbirthname}{in}{in}
\bilingual{\coursename}{Studiengang:}{Course:}
\bilingual{\disciplinename}{Studienrichtung:}{Discipline:}
\bilingual{\matriculationnumbername}{Matrikelnummer:}{Matriculation number:}
\bilingual{\matriculationyearname}{Immatrikulationsjahr}{Matriculation year:}
\bilingual{\graduationname}{zur Erlangung des akademischen Grades}{to achieve the academic degree}
\bilingual{\refereename}{Gutachter}{Referee}
\bilingual{\advisorname}{Fachreferent}{Advisor}
\bilingual{\supervisorname}{Betreuer}{Supervisor}
\bilingual{\professorname}{Betreuender Hochschullehrer}{Supervising professor}
\bilingual{\datename}{Eingereicht am:}{Submitted on:}
%% Thesis Typen (inspiriert von dem Paket tudscr):
\bilingual{\dissertationname}{Dissertation}{Dissertation} % diss, doctoral, phd 
\bilingual{\diplomathesisname}{Diplomarbeit}{Diploma Thesis} % diploma 
\bilingual{\masterthesisname}{Master-Arbeit}{Master Thesis} % master 
\bilingual{\bachelorthesisname}{Bachelor-Arbeit}{Bachelor Thesis} % bachelor 
\bilingual{\studentthesisname}{Studienarbeit}{Student Thesis} % student 
\bilingual{\studentresearchname}{Großer Beleg}{Student Research Project} % evidence 
\bilingual{\projectpapername}{Projektarbeit}{Project Paper} % project
\bilingual{\seminarpapername}{Seminararbeit}{Seminar Paper} % seminar
\bilingual{\termpapername}{Hausarbeit}{Term Paper} % term
\bilingual{\researchname}{Forschungsbericht}{Research Report} % research 
\bilingual{\logname}{Protokoll}{Log} % log 
\bilingual{\reportname}{Bericht}{Report} % report
\bilingual{\internshipname}{Praktikumsbericht}{Internship Report} % internship 
\bilingual{\lecturename}{Vorlesungsmitschrift}{Lecture Notes} % lecture 
\bilingual{\lectureauthorname}{Mitschrift von}{Notes by} % lecture  
\bilingual{\lessonname}{\"Ubungsmitschrift}{Lesson Notes} % lecture  
\bilingual{\practname}{Praktikumsmitschrift}{Pracital Course Notes} % lecture 
%%% Variablen für die Titelseite (inspiriert von dem Paket tudscr):
\renewcommand{\subject}[1]{\checkThesis{#1}{\@subject}}
\newcommand{\faculty}[1]{\gdef\@faculty{#1}}
\newcommand{\department}[1]{\gdef\@department{#1}}
\newcommand{\institute}[1]{\gdef\@institute{#1}}
\newcommand{\chair}[1]{\gdef\@chair{#1}}
\newcommand{\authormore}[1]{\gdef\@authormore{#1}}
\newcommand{\dateofbirth}[1]{\gdef\@dateofbirth{#1}}
\newcommand{\placeofbirth}[1]{\gdef\@placeofbirth{#1}}
\newcommand{\course}[1]{\gdef\@course{#1}}
\newcommand{\discipline}[1]{\gdef\@discipline{#1}}
\newcommand{\matriculationnumber}[1]{\gdef\@matriculationnumber{#1}}
\newcommand{\matriculationyear}[1]{\gdef\@matriculationyear{#1}}
\newcommand{\thesis}[1]{\checkThesis{#1}{\@thesis}}
\newcommand{\graduation}[2][]{\gdef\@graduation{#2}\gdef\@graduationshort{#1}}
\newcommand{\supervisor}[1]{\gdef\@supervisor{#1}}
\newcommand{\referee}[1]{\gdef\@referee{#1}}
\newcommand{\advisor}[1]{\gdef\@advisor{#1}}
\newcommand{\professor}[1]{\gdef\@professor{#1}}
%%% Titelseite (inspiriert von dem Paket tudscr):
\renewcommand*{\maketitle}{%
\begin{titlepage}%
\begin{flushleft}
\renewcommand{\thefootnote}{$\star$}%
%%% Kopfzeile HTW-Style:
\@htwheader\\[13mm]
%%% Die Zeilen für titlehead, subject, title, subtitle und author sind immer vorhanden (spacing)
{\noindent\@titlehead{}\strut}\\[3mm]%
{\Large\@subject{}\strut}\\[8mm]%
{\Huge\bfseries%
\ifdefined\@nouppercase\@title{}\else\MakeUppercase{\@title{}}\fi%
\strut}\\
{\Large\bfseries%
\ifdefined\@nouppercase\@subtitle{}\else\MakeUppercase{\@subtitle{}}\fi%
\strut}\\[11mm]
\ifdefined\@islecture\strut\else{\Large\@author{}\strut}\fi % Autor wird bei Vorlesungsmitschrift unten angezeigt
%%% Weitere Felder werden nur angezeigt, wenn sie definiert wurden
\ifdefined\@authormore\ifdefined\@islecture\else\\\fi{\@authormore{}\strut}\fi
\ifdefined\@dateofbirth\\\dateofbirthname{} \@dateofbirth{} \ifdefined\@placeofbirth\placeofbirthname{} \@placeofbirth{}\strut\fi\fi
\ifdefined\@course\\\coursename{} \@course{}\strut\fi
\ifdefined\@discipline\\\disciplinename{} \@discipline{}\strut\fi
\ifdefined\@matriculationnumber\\\matriculationnumbername{} \@matriculationnumber{}\strut\fi
\ifdefined\@matriculationyear\\\matriculationyearname{} \@matriculationyear{}\strut\fi
\ifdefined\@thesis\\[10mm]{\LARGE\bfseries
\ifdefined\@nouppercase\@thesis{}\else\MakeUppercase{\@thesis{}}\fi
\strut}
\ifdefined\@graduation\\[2mm]{\graduationname{}\strut}
\\[5mm]{\LARGE\bfseries
\ifdefined\@nouppercase\relax\else\MakeUppercase{\fi%
\@graduation{} \ifdefined\@graduationshort(\@graduationshort{})\fi
\ifdefined\@nouppercase\relax\else}\fi\strut}\fi\fi
\vfill\noindent
\ifdefined\@islecture\\[1mm]{\lectureauthorname{}\strut}\\[.7mm]{\Large\@author{}\strut}\fi
\ifdefined\@referee{\refereename{}\strut}\\[.7mm]{\Large\@referee{}\strut}\fi
\ifdefined\@advisor\\[1mm]{\advisorname{}\strut}\\[.7mm]{\Large\@advisor{}\strut}\fi
\ifdefined\@supervisor\\[1mm]{\supervisorname{}\strut}\\[.7mm]{\Large\@supervisor{}\strut}\fi
\ifdefined\@professor\\[1mm]{\professorname{}\strut}\\[.7mm]{\Large\@professor{}\strut}\fi
\ifdefempty{\@date}{}{\ifdefined\@simpledate\\[6mm]\relax\else\\[3mm]\datename{}~\fi\@date{}\strut}
\ifdefempty{\@publishers}{}{\\[4.5mm]{\Large\@publishers{}\strut}}%
\renewcommand\footnoterule{}
\@thanks\global\let\@thanks\@empty
\end{flushleft}
\end{titlepage}
\setcounter{footnote}{0}
}
\else%
\PackageError{htwcd}{htwcd needs a KOMA class!}%
\fi%
\makeatother
%%% === Ende Paket htwscr ===
%-- Nummerierung bis Subsubsection für Report
\setcounter{secnumdepth}{3}	% zählt auch subsubsection
\setcounter{tocdepth}{3}	% Inhaltsverzeichnis bis in subsubsection

% LuaLaTeX Unterscheidung
\usepackage{ifluatex}

% Input inkl. Umlaute, Silbentrennung
\ifluatex
	\usepackage{fontspec}
\else
	\usepackage[utf8]{inputenc}
	\usepackage[T1]{fontenc}
\fi
\usepackage[english,ngerman]{babel}
\usepackage{csquotes}	% Anführungszeichen
\usepackage{marvosym}
\usepackage{eurosym}

% Style-Aufhübschung
\usepackage{soul, color}	% Kapitälchen, Unterstrichen, Durchgestrichen usw. im Text
%\usepackage{titleref}
\usepackage[breakwords, fit]{truncate}	% Abschneiden von Sätzen
\renewcommand{\TruncateMarker}{\,…}

% Mathe usw.
\usepackage{amssymb}
\usepackage{amsthm}
\usepackage[fleqn,intlimits]{amsmath}	% fleqn: align-Umgebung rechtsbündig; intlimits: Integralgrenzen immer ober-/unterhalb
%\usepackage{mathtools} % u.a. schönere underbraces
\usepackage{xcolor}
\usepackage{esint}	% Schönere Integrale, \oiint vorhanden
\everymath=\expandafter{\the\everymath\displaystyle}	% Mathe Inhalte werden weniger verkleinert
\usepackage{wasysym}	% mehr Symbole, bspw \lightning
%\renewcommand{\int}{\int\limits}
%\usepackage{xfrac}	% mehr fracs: sfrac{}{}
\let\oldemptyset\emptyset	% schöneres emptyset
\let\emptyset\varnothing
%\RequirePackage{mathabx}	% mehr Symbole
\mathchardef\mhyphen="2D	% Hyphen in Math
\usepackage{mathtools}		% zum Beispiel {pmatrix}[r] für Formatierung der Matrix

% tikz usw.
\usepackage{tikz}
\usepackage{pgfplots}
\pgfplotsset{compat=1.11}	% Umgeht Fehlermeldung
\usetikzlibrary{graphs}
%\usetikzlibrary{through}	% ???
\usetikzlibrary{arrows}
\usetikzlibrary{arrows.meta}	% Pfeile verändern / vergrößern: \draw[-{>[scale=1.5]}] (-3,5) -> (-3,3);
\usetikzlibrary{automata,positioning} % Zeilenumbruch im Node node[align=center] {Text\\nächste Zeile} automata für Graphen
\usetikzlibrary{matrix}
\usetikzlibrary{patterns}	% Schraffierte Füllung
\usetikzlibrary{shapes.geometric}	% Polygon usw.
\tikzstyle{reverseclip}=[insert path={	% Inverser Clip \clip
	(current page.north east) --
	(current page.south east) --
	(current page.south west) --
	(current page.north west) --
	(current page.north east)}
% Nutzen: 
%\begin{tikzpicture}[remember picture]
%\begin{scope}
%\begin{pgfinterruptboundingbox}
%\draw [clip] DIE FLÄCHE, IN DER OBJEKT NICHT ERSCHEINEN SOLL [reverseclip];
%\end{pgfinterruptboundingbox}
%\draw DAS OBJEKT;
%\end{scope}
%\end{tikzpicture}
]	% Achtung: dafür muss doppelt kompliert werden!
\usepackage{graphpap}	% Grid für Graphen
\tikzset{every state/.style={inner sep=2pt, minimum size=2em}}
\usetikzlibrary{mindmap, backgrounds}
%\usepackage{tikz-uml}	% braucht Dateien: http://perso.ensta-paristech.fr/~kielbasi/tikzuml/

% Tabular
\usepackage{longtable}	% Große Tabellen über mehrere Seiten
\usepackage{multirow}	% Multirow/-column: \multirow{2[Anzahl der Zeilen]}{*[Format]}{Test[Inhalt]} oder \multicolumn{7[Anzahl der Reihen]}{|c|[Format]}{Test2[Inhalt]}
\renewcommand{\arraystretch}{1.3} % Tabellenlinien nicht zu dicht
\usepackage{colortbl}
\arrayrulecolor{gray}	% heller Tabellenlinien
\usepackage{array}	% für folgende 3 Zeilen (für Spalten fester breite mit entsprechender Ausrichtung):
\newcolumntype{L}[1]{>{\raggedright\let\newline\\\arraybackslash\hspace{0pt}}m{\dimexpr#1\columnwidth-2\tabcolsep-1.5\arrayrulewidth}}
\newcolumntype{C}[1]{>{\centering\let\newline\\\arraybackslash\hspace{0pt}}m{\dimexpr#1\columnwidth-2\tabcolsep-1.5\arrayrulewidth}}
\newcolumntype{R}[1]{>{\raggedleft\let\newline\\\arraybackslash\hspace{0pt}}m{\dimexpr#1\columnwidth-2\tabcolsep-1.5\arrayrulewidth}}
\usepackage{booktabs}	% Kompatibilität für array und \hrule
\usepackage{caption}	% Um auch unbeschriftete Captions mit \caption* zu machen

% Nützliches
\usepackage{verbatim}	% u.a. zum auskommentieren via \begin{comment} \end{comment}
\usepackage{tabto}	% Tabs: /tab zum nächsten Tab oder /tabto{.5 \CurrentLineWidth} zur Stelle in der Linie
\NumTabs{6}	% Anzahl von Tabs pro Zeile zum springen
\usepackage{listings} % Source-Code mit Tabs
\usepackage{lstautogobble} 
\usepackage{enumitem}	% Anpassung der enumerates
%\setlist[enumerate,1]{label=(\arabic*)}	% global andere Enum-Items
\renewcommand{\labelitemiii}{$\scriptscriptstyle ^\blacklozenge$} % global andere 3. Item-Aufzählungszeichen
\usepackage{letltxmacro} % neue Definiton von Grundbefehlen
% Nutzen:
%\LetLtxMacro{\oldemph}{\emph}
%\renewcommand{\emph}[1]{\oldemph{#1}}
\RequirePackage{xpatch}	% ua. Konkatenieren von Strings/Variablen (etoolbox)
\usepackage{xstring}	% String Operationen
\usepackage{minibox}	% Minibox anstatt \fbox{} für Boxen mit (manuellen) Zeilenumbruch
\usepackage{framed}		% \begin{framed} statt minibox/fbox mit automatischem Zeilenumbruch
\usepackage{microtype}	% Verfeinerte Hyphenation
\usepackage{import}		
\usepackage{standalone}	% Zusammen mit import: Einfügen von separaten Latex-Dateien (die schon header haben)

% Einrichtung von lst
%\renewcommand*{\ttdefault}{qcr}
\lstset{
basicstyle=\ttfamily, 
%basicstyle=\footnotesize\sffamily\color{black},
%mathescape=true, 
%escapeinside=^^, 
autogobble, 
tabsize=2,
frame=single,
rulecolor=\color{lightgray},
numbers=left,
numbersep=5pt,
numberstyle=\tiny\color{gray},
commentstyle=\color{gray},
keywordstyle=\color{htwblue},
stringstyle=\color{htworange},
morecomment=[l][\color{magenta}]{\#}
showspaces=false,
showstringspaces=false,
breaklines=true,
literate=%
    {Ö}{{\"O}}1
    {Ä}{{\"A}}1
    {Ü}{{\"U}}1
    {ß}{{\ss}}1
    {ü}{{\"u}}1
    {ä}{{\"a}}1
    {ö}{{\"o}}1
    {„}{{\glqq}}1
    {“}{{\grqq}}1
    {~}{{\textasciitilde}}1
}
\usepackage{scrhack} % Fehler umgehen
\def\ContinueLineNumber{\lstset{firstnumber=last}} % vor lstlisting. Zum wechsel zum nicht-kontinuierlichen muss wieder \StartLineAt1 eingegeben werden
\def\StartLineAt#1{\lstset{firstnumber=#1}} % vor lstlisting \StartLineAt30 eingeben, um bei Zeile 30 zu starten
\let\numberLineAt\StartLineAt

% BibTeX
\usepackage[
backend=bibtex,
bibencoding=ascii,
style=authoryear, citestyle=authoryear-ibid,
%doi=false,
%isbn=false,
%url=false
]{biblatex}	% BibTeX
\usepackage{makeidx}
%\makeglossary
%\makeindex
\usepackage{url}	% Darstellung von URLs

% Grafiken
\usepackage{graphicx}
\usepackage{epstopdf}	% eps-Vektorgrafiken einfügen
\usepackage{transparent}	% transparent nutzen: {\transparent{0.4} ...}
%\epstopdfsetup{outdir=\customDir}
% Prüft, ob Grafik existiert (mit \ifvalidimage{}{}) [Quelle: https://tex.stackexchange.com/a/99176]:
\makeatletter
\newif\ifgraphicexist
\catcode`\*=11
\newcommand\ifvalidimage[1]{%
    \begingroup
    \global\graphicexisttrue
    \let\input@path\Ginput@path
    \filename@parse{#1}%
    \ifx\filename@ext\relax
    \@for\Gin@temp:=\Gin@extensions\do{%
        \ifx\Gin@ext\relax
        \Gin@getbase\Gin@temp
        \fi}%
    \else
    \Gin@getbase{\Gin@sepdefault\filename@ext}%
    \ifx\Gin@ext\relax
    \global\graphicexistfalse
    \def\Gin@base{\filename@area\filename@base}%
    \edef\Gin@ext{\Gin@sepdefault\filename@ext}%
    \fi
    \fi
    \ifx\Gin@ext\relax
    \global\graphicexistfalse
    \else 
    \@ifundefined{Gin@rule@\Gin@ext}%
    {\global\graphicexistfalse}%
    {}%
    \fi  
    \ifx\Gin@ext\relax 
    \gdef\imageextension{unknown}%
    \else
    \xdef\imageextension{\Gin@ext}%
    \fi 
    \endgroup 
    \ifgraphicexist
    \expandafter \@firstoftwo
    \else
    \expandafter \@secondoftwo
    \fi 
} 
\catcode`\*=12
\makeatother
\LetLtxMacro{\forceincludegraphics}{\includegraphics}	% neuer Befehl für includegraphics
\renewcommand{\includegraphics}[2][]{	% altes includegraphics neu definieren, damit es auch nicht vorhandene einfügt
\ifvalidimage{#2}{
\forceincludegraphics[#1]{#2}
}{
\message{Achtung: Grafik wurde nicht gefunden: '#2'}
\minibox[frame]{
\textbf{\StrSubstitute{#2}{_}{\_}}%  \ifthenelse{\isempty{#1}}{}{\\\textit{#1}}
}}}

% pdf-Setup
\usepackage{pdfpages}
\usepackage[bookmarks,%
bookmarksopen=false,% Klappt die Bookmarks in Acrobat aus
colorlinks=true,%
linkcolor=black,%
citecolor=black,%
urlcolor=htworange,%
]{hyperref}

% Nach Dokumentenbeginn ausführen:
\makeatletter
\AtBeginDocument{	
	% pdf Informationen
	\hypersetup{
		pdfauthor={\@author},
		pdftitle={\@title},
		pdfproducer={LaTeX Layout erstellt von Falk-Jonatan Strube}
	}
	% signierte Fußzeile
	%\providecommand{\ofootname}{Mitschrift von}
	\providecommand{\pathtomaster}{../}
	\ofoot{%
	%\ifdefined\ofoottext
	%\else
	%\begin{minipage}{40mm}
	%\begin{flushright}
	%\textcolor{gray}{
	%\ofootname\\
	%Falk-Jonatan Strube\vspace*{-1mm}}
	%\end{flushright}
	%\end{minipage}
	%\fi%
	%\hspace*{1mm}%
	\begin{minipage}{10 mm}
	\includegraphics[height=9mm]{\pathtomaster/FJS-logo_gray}
	\end{minipage}%
	}
}
\makeatother



%%% ==================================================
%%% ===================== MACROS =====================
%%% ==================================================
% Horizontale Linie:
\newcommand{\HRule}[1][\medskipamount]{\par
  \vspace*{\dimexpr-\parskip-\baselineskip+#1}
  \noindent\rule[0.2ex]{\linewidth}{0.2mm}\par
  \vspace*{\dimexpr-\parskip-.5\baselineskip+#1}}
% Gestrichelte horizontale Linie:
\RequirePackage{dashrule}
\newcommand{\HDRule}[1][\medskipamount]{\par
  \vspace*{\dimexpr-\parskip-\baselineskip+#1}
  \noindent\hdashrule[0.2ex]{\linewidth}{0.2mm}{1mm} \par
  \vspace*{\dimexpr-\parskip-.5\baselineskip+#1}}
% Mathe in Anführungszeichen:
\newsavebox{\mathbox}\newsavebox{\mathquote}
\makeatletter
\newcommand{\mq}[1]{% \mathquotes{<stuff>}
  \savebox{\mathquote}{\text{"}}% Save quotes
  \savebox{\mathbox}{$\displaystyle #1$}% Save <stuff>
  \raisebox{\dimexpr\ht\mathbox-\ht\mathquote\relax}{"}#1\raisebox{\dimexpr\ht\mathbox-\ht\mathquote\relax}{''}
}
\makeatother

% Paragraph mit Zähler (Section-Weise)
\newcounter{cparagraphC}
\newcommand{\cparagraph}[1]{
\stepcounter{cparagraphC}
\paragraph{\thesection{}-\thecparagraphC{} #1}
%\addcontentsline{toc}{subsubsection}{\thesection{}-\thecparagraphC{} #1}
\label{\thesection-\thecparagraphC}
}
\makeatletter
\@addtoreset{cparagraphC}{section}
\makeatother


% (Vorlesungs-)Folien einbinden:
\makeatletter
% Folien-Pfad (inkl. Dateiname ohne Endung und ggf. ohne Nummerierung)
\newcommand{\@slidePath}{}
\newcommand{\slidePath}[1]{\renewcommand{\@slidePath}{#1}}
% Folien Eigenschaften
\newcommand{\@slideScale}{0.5}
\newcommand{\slideScale}[1]{\renewcommand{\@slideScale}{#1}}
\AtBeginDocument{
% Folien von einer Datei skaliert
\newcommand{\slide}[2][\@slideScale]{\slides[#1]{}{#2}}
\newcommand{\slideTrim}[6][\@slideScale]{\slides[#1 , clip,  trim = #5cm #4cm #6cm #3cm]{}{#2}}
% Folien von mehreren nummerierten Dateien skaliert
\newcommand{\slides}[3][\@slideScale]{\begin{center}
\includegraphics[page=#3, scale=#1]{\@slidePath #2.pdf}
\end{center}}
}
\makeatother

% \emph{} anders definieren
\makeatletter
\DeclareRobustCommand{\em}{%
  \@nomath\em \if b\expandafter\@car\f@series\@nil
  \normalfont \else \scshape \fi}
\makeatother

% unwichtiges
\newcommand{\unimptnt}[1]{{\transparent{0.5}#1}}

% alph. enumerate
\newenvironment{anumerate}{\begin{enumerate}[label=(\alph*)]}{\end{enumerate}} % Alphabetische Aufzählung

% Hanging parameters
\newcommand{\hangpara}[1]{\par\noindent\hangindent+2em\hangafter=1 #1\par\noindent}

%% EINFACHE BEFEHLE

% Abkürzungen Mathe
\newcommand{\EE}{\mathbb{E}}
\newcommand{\QQ}{\mathbb{Q}}
\newcommand{\RR}{\mathbb{R}}
\newcommand{\CC}{\mathbb{C}}
\newcommand{\NN}{\mathbb{N}}
\newcommand{\ZZ}{\mathbb{Z}}
\newcommand{\PP}{\mathbb{P}}
\renewcommand{\SS}{\mathbb{S}}
\newcommand{\cA}{\mathcal{A}}
\newcommand{\cB}{\mathcal{B}}
\newcommand{\cC}{\mathcal{C}}
\newcommand{\cD}{\mathcal{D}}
\newcommand{\cE}{\mathcal{E}}
\newcommand{\cF}{\mathcal{F}}
\newcommand{\cG}{\mathcal{G}}
\newcommand{\cH}{\mathcal{H}}
\newcommand{\cI}{\mathcal{I}}
\newcommand{\cJ}{\mathcal{J}}
\newcommand{\cM}{\mathcal{M}}
\newcommand{\cN}{\mathcal{N}}
\newcommand{\cP}{\mathcal{P}}
\newcommand{\cR}{\mathcal{R}}
\newcommand{\cS}{\mathcal{S}}
\newcommand{\cZ}{\mathcal{Z}}
\newcommand{\cL}{\mathcal{L}}
\newcommand{\cT}{\mathcal{T}}
\newcommand{\cU}{\mathcal{U}}
\newcommand{\cX}{\mathcal{X}}
\newcommand{\cV}{\mathcal{V}}
\renewcommand{\phi}{\varphi}
\renewcommand{\epsilon}{\varepsilon}
\renewcommand{\theta}{\vartheta}

% Verschiedene als Mathe-Operatoren
\DeclareMathOperator{\arccot}{arccot}
\DeclareMathOperator{\arccosh}{arccosh}
\DeclareMathOperator{\arcsinh}{arcsinh}
\DeclareMathOperator{\arctanh}{arctanh}
\DeclareMathOperator{\arccoth}{arccoth} 
\DeclareMathOperator{\var}{Var} % Varianz 
\DeclareMathOperator{\cov}{Cov} % Co-Varianz 

% Farbdefinitionen
\definecolor{red}{RGB}{180,0,0}
\definecolor{green}{RGB}{75,160,0}
\definecolor{blue}{RGB}{0,75,200}
\definecolor{orange}{RGB}{255,128,0}
\definecolor{yellow}{RGB}{255,245,0}
\definecolor{purple}{RGB}{75,0,160}
\definecolor{cyan}{RGB}{0,160,160}
\definecolor{brown}{RGB}{120,60,10}

\definecolor{itteny}{RGB}{244,229,0}
\definecolor{ittenyo}{RGB}{253,198,11}
\definecolor{itteno}{RGB}{241,142,28}
\definecolor{ittenor}{RGB}{234,98,31}
\definecolor{ittenr}{RGB}{227,35,34}
\definecolor{ittenrp}{RGB}{196,3,125}
\definecolor{ittenp}{RGB}{109,57,139}
\definecolor{ittenpb}{RGB}{68,78,153}
\definecolor{ittenb}{RGB}{42,113,176}
\definecolor{ittenbg}{RGB}{6,150,187}
\definecolor{itteng}{RGB}{0,142,91}
\definecolor{ittengy}{RGB}{140,187,38}

% Textfarbe ändern
\newcommand{\tred}[1]{\textcolor{red}{#1}}
\newcommand{\tgreen}[1]{\textcolor{green}{#1}}
\newcommand{\tblue}[1]{\textcolor{blue}{#1}}
\newcommand{\torange}[1]{\textcolor{orange}{#1}}
\newcommand{\tyellow}[1]{\textcolor{yellow}{#1}}
\newcommand{\tpurple}[1]{\textcolor{purple}{#1}}
\newcommand{\tcyan}[1]{\textcolor{cyan}{#1}}
\newcommand{\tbrown}[1]{\textcolor{brown}{#1}}

% Umstellen der Tabellen Definition
\newcommand{\mpb}[1][.3]{\begin{minipage}{#1\textwidth}\vspace*{3pt}}
\newcommand{\mpe}{\vspace*{3pt}\end{minipage}}

\newcommand{\resultul}[1]{\underline{\underline{#1}}}
\newcommand{\parskp}{$ $\\}	% new line after paragraph
\newcommand{\parnl}{\unskip\parfillskip 0pt \par} % justified new line (brace line!)

% mathematische Abkürzungen
\newcommand{\corr}{\;\widehat{=}\;}
\newcommand{\mdeg}{^{\circ}}
\newcommand{\nok}[2]{\binom{#1}{#2}}	% n über k BESSER: \binom{n}{k}
\newcommand{\mtr}[1]{\begin{pmatrix}#1\end{pmatrix}}	% Matrix
\newcommand{\dtr}[1]{\begin{vmatrix}#1\end{vmatrix}}	% Determinante (Betragsmatrix)
\LetLtxMacro{\originalVec}{\vec}
\renewcommand{\vec}[1]{\underline{#1}}	% Vektorschreibweise
\newcommand{\imptnt}[1]{\colorbox{red!30}{#1}}	% Wichtiges
\newcommand{\intd}[1]{\,\mathrm{d}#1}
\newcommand{\diffd}[1]{\mathrm{d}#1}
% für Module-Rechnung: \pmod{}
\newcommand{\unit}[1]{\,\mathrm{#1}}
%\LetLtxMacro{\ntilde}{\tilde}
%\renewcommand{\tilde}{\widetilde}

% weiter Abkürzungen
\newcommand{\gdw}{genau dann wenn}
\newcommand{\lecdate}[1]{\marginpar{\small\textcolor{gray}{Vorlesung\newline #1}}}

\thesis{Entwicklerdokumentation}
\date{}

\makeatletter
\gdef\@islecture{\relax}
\makeatother

\faculty{Fakultät Informatik/Mathematik}
\chair{Lehrstuhl für Softwaretechnologie}
\subject{Software Engineering 2}
\title{LED-Leuchtturm}
%\subtitle{Entwicklung eines LED-Leuchtturms zur Anzeige des aktuellen Maschinenzustands}
\bilingual{\lectureauthorname}{Gruppenmitglieder}{Group members}
\author{
\texorpdfstring{Ragnar Luga		\tabto{.35\linewidth} (s75594)\\
Raphael Pour									\tabto{.35\linewidth} (s74020)\\
Torsten Polczyk								\tabto{.35\linewidth} (s74097)\\
Simon Retsch									\tabto{.35\linewidth} (s74077)\\
Eric Schmidtgen								\tabto{.35\linewidth} (s73952)\\
Ben Schönherr								\tabto{.35\linewidth} (s74083)\\
Falk-Jonatan Strube						\tabto{.35\linewidth} (s74053)\\
Jürgen Tomaszewski						\tabto{.35\linewidth} (s76451)\\\vspace*{-1em}}{
Ragnar Luga,
Raphael Pour,
Torsten Polczyk,
Simon Retsch,
Erik Schmidtgen,
Ben Schoenherr,
Falk-Jonatan Strube,
Jürgen Tomaszewski
}
}
\bilingual{\supervisorname}{Auftraggeber}{Client}
\supervisor{Dipl.-Inf. Robert Ringel}
\professor{Prof. Dr. Anna Sabine Hauptmann}

\let\tableofcontentsORIG\tableofcontents
\renewcommand\tableofcontents{\tableofcontentsORIG\clearpage}

\begin{document}
\maketitle
\tableofcontents
\renewcommand{\maketitle}{}
\renewcommand{\tableofcontents}{}
% Anmerkung: Entwicklerdokumentation = Programmiererdokumentation (Pflichtenheft /DA2/)

\section*{Einleitung}

In der Entwicklerdokumentation wird die Entwicklung des Softwaresystems \texttt{ledleuchtturm} dargestellt. In den Abschnitten \nameref{grobentw} und \nameref{feinentw} sind die Vorüberlegungen zu der Software zu finden. Diese Abschnitte erläutern, auf welche Themen geachtet wurde und wie ein Plan für das finale Produkt entstand.

Im Abschnitt \nameref{strukt} ist dargestellt, wie das Softwaresystem, bzw. dessen Code, im Vergleich zum Entwurf organisiert ist. Der Abschnitt \nameref{epydoc} stellt schließlich die vollständige Dokumentation des Codes dar.

%Diese Dokumentation soll dabei helfen einen Einblick in das Softwaresystem zu bekommen.

\section{Grobentwurf}
\label{grobentw}
% LaTeX master Datei(en) zusammengestellt von Falk-Jonatan Strube zur Nutzung an der Hochschule für Technik und Wirtschaft Dresden: https://github.com/genericFJS/htw
\documentclass{scrartcl}
\gdef\pathtomaster{../_LaTeX_master}
%% HTW Dresden Corporate Design
%% Copyright (C) Falk-Jonatan Strube <falkjonatanstrube@gmail.com>, 2017
%%% === Anfang Paket htwscr ===
\makeatletter
\ifcsname KOMAoption\endcsname%
\KOMAoption{footheight}{11mm}	% mehr Platz für Wasserzeichen o.ä. in der Fußzeile
\recalctypearea	% s.o.
\RequirePackage{geometry}
\geometry{top=25mm, bottom=25mm, headsep=10mm, inner=25mm, outer=25mm, footskip=15mm, marginparsep=2mm}
\RequirePackage{helvet}	% Corporate Design schreibt Arial/Helvetica vor
\renewcommand{\familydefault}{\sfdefault}	% s.o.
\RequirePackage{graphicx}	% für Grafiken
\RequirePackage{xcolor}		% für Farben
\definecolor{htworange}{RGB}{249,155,28}	% Corporate Design orange CMYK(0/45/100/0)
\definecolor{htwblue}{RGB}{0,116,188}			% Corporate Design blau CMYK(90/50/0/0)
\definecolor{htwgrey}{RGB}{128,128,128}		% Corporate Design gray 60% Schwarz [Annahme: CMYK(0,0,0,60)]
\RequirePackage{ifthen,xifthen}	% für Bedingungen/Optionen
\RequirePackage{etoolbox}				% s.o.

%%% Format des HTW-Headers (wird hier als Texteelement verwenden) nach dem Style der PowerPoint-Vorlagen. Weiterhin inspiriert durch das Paket tudscr (Corporate Design der TU Dresden)
%%% Dateipfad des Logos abhängig von der Verwendung als Paket/Input
\ifdefined\@isHTWCDpackage%
	\gdef\@pathtoimage{HTW}%
\else%
	\ifdefined\pathtomaster%
		\gdef\@pathtoimage{\pathtomaster /HTW}%
	\else 
		\gdef\@pathtoimage{HTW}%
	\fi%
\fi%
\newcommand{\pathtologo}[1]{\gdef\@pathtoimage{#1}}
%%% Berechnung der Größe des Logos
\newlength{\htwLogoHeight}
\setlength{\htwLogoHeight}{12mm}	% Mindesthöhe des Logos laut Corporate Design: 11mm
\newlength{\@htwLogoWidth}
\AtBeginDocument{
\def\printLogo{\includegraphics[height=\htwLogoHeight]{\@pathtoimage}}
\setlength\@htwLogoWidth{\widthof{\printLogo}}
}
%%% Definition des Header
\newcommand{\@htwheader}{{\noindent%
%%% Header ist abhängig von dem Vorhandensein von department, institute und chair
\newcommand{\hasDepInsCh}{\gdef\@hasDepInsCh{}}%
\ifdefined\@department\hasDepInsCh\fi%
\ifdefined\@institute\hasDepInsCh\fi%
\ifdefined\@chair\hasDepInsCh\fi%
\begin{minipage}[b][\htwLogoHeight]{\dimexpr \textwidth - \@htwLogoWidth - 2mm \relax}% 2mm Abstand zwischen Logo und Linie
{\small{\bfseries%
\ifdefined\@faculty\@faculty{}\fi} %
\strut}%
{\color{htworange}\hrule height 1pt}%
\vspace{.1mm}%
\end{minipage}
\hfill%
\begin{minipage}[b][\htwLogoHeight]{\@htwLogoWidth}
\includegraphics[height=\htwLogoHeight]{\@pathtoimage}
\end{minipage}
\ifdefined\@hasDepInsCh%
\\%
\begin{minipage}{\textwidth}
\small%
\ifdefined\@department\@department{}\fi%
\ifdefined\@institute\ifdefined\@department,~\fi \@institute\fi%
\ifdefined\@chair\ifdefined\@department,~\else\ifdefined\@institute,~\fi\fi\@chair\fi\strut
%{\color{htworange}\hrule height 0.5pt}
\end{minipage}
\fi
}}

%%% Definition der Fußzeilen
\usepackage{scrlayer-scrpage}
\@ifclassloaded{scrreprt}{%
\renewcommand*{\partpagestyle}{scrheadings}
\renewcommand*{\chapterpagestyle}{scrheadings}%
}{}
\pagestyle{scrheadings}
\cfoot{\vfill\pagemark}
\chead{}

%%% Einstellungen
\newcommand{\simpledate}{\gdef\@simpledate{}}		% Datum nicht mit "Eingereicht am:" Vorsatz
\newcommand{\nouppercase}{\gdef\@nouppercase{}}	% Titel/Parts/Chapters usw. nicht in Großbuchstaben

%%% Einiges erst zu Beginn des Dokumentes ausführen, um Einstellungen zu berücksichtigen
\AtBeginDocument{
\renewcommand\SS{SS}	% fix big ß
%%% Sektionen usw. ggf. in Großbuchstaben:
\providecommand{\chapterlinesformat}[3]{}	% compatibility
\providecommand{\chapterlineswithprefixformat}[3]{}	% compatibility
\providecommand{\sectionlinesformat}[4]{}	% compatibility
\ifdefined\@nouppercase\relax\else
\renewcommand\sectionlinesformat[4]{\@hangfrom{\hskip#2 #3}{\MakeUppercase{#4}}}
\@ifclassloaded{scrreprt}{%
\renewcommand\chapterlinesformat[3]{\@hangfrom{#2}{\MakeUppercase{#3}}}
\renewcommand\chapterlineswithprefixformat[3]{\MakeUppercase{#2#3}}%
}{}
\fi
%%% Anpassung der Part-Seite
\renewcommand*{\partheadstartvskip}{\@htwheader\vspace*{30mm}}
\renewcommand*{\raggedpart}{\raggedright}
\renewcommand*{\partformat}{\huge\ifdefined\@nouppercase \partname~\thepart\else\MakeUppercase{\partname~\thepart}\fi.}
\ifdefined\@nouppercase\setkomafont{part}{\Huge}\else\setkomafont{part}{\Huge\MakeUppercase}\fi
}

\newcommand{\bilingual}[3]{\ifdefined #1 \let #1\undefined\fi\newcommand{#1}{\iflanguage{english}{#3}{#2}}}	% Übersetzung von einigen Variablen in mehrere Sprachen

\newcommand{\checkThesis}[2]{\ifthenelse{\equal{#1}{diss}\OR\equal{#1}{doctoral}\OR\equal{#1}{phd}}{
\gdef#2{\dissertationname}}{
\ifthenelse{\equal{#1}{diploma}}{
\gdef#2{\diplomathesisname}}{
\ifthenelse{\equal{#1}{master}}{
\gdef#2{\masterthesisname}}{
\ifthenelse{\equal{#1}{bachelor}}{
\gdef#2{\bachelorthesisname}}{
\ifthenelse{\equal{#1}{student}}{
\gdef#2{\studentthesisname}}{
\ifthenelse{\equal{#1}{evidence}}{
\gdef#2{\studentresearchname}}{
\ifthenelse{\equal{#1}{project}}{
\gdef#2{\projectpapername}}{
\ifthenelse{\equal{#1}{seminar}}{
\gdef#2{\seminarpapername}}{
\ifthenelse{\equal{#1}{term}}{
\gdef#2{\termpapername}}{
\ifthenelse{\equal{#1}{research}}{
\gdef#2{\researchname}}{
\ifthenelse{\equal{#1}{log}}{
\gdef#2{\logname}}{
\ifthenelse{\equal{#1}{report}}{
\gdef#2{\reportname}}{
\ifthenelse{\equal{#1}{internship}}{
\gdef#2{\internshipname}}{
\ifthenelse{\equal{#1}{lecture}\OR\equal{#1}{lesson}\OR\equal{#1}{pract}}{
\bilingual{\professorname}{Vorlesung von}{Lecture by}
\simpledate
\gdef\@islecture{\relax}
\ifthenelse{\equal{#1}{lecture}}{\gdef#2{\lecturename}}{
\ifthenelse{\equal{#1}{lesson}}{\gdef#2{\lessonname}}{
\gdef#2{\practname}
}}}{
\gdef#2{#1}}}}}}}}}}}}}}}}

%%% Textfragmente der Titelseite in deutsch und englisch
\bilingual{\dateofbirthname}{Geboren am:}{Born on:}
\bilingual{\placeofbirthname}{in}{in}
\bilingual{\coursename}{Studiengang:}{Course:}
\bilingual{\disciplinename}{Studienrichtung:}{Discipline:}
\bilingual{\matriculationnumbername}{Matrikelnummer:}{Matriculation number:}
\bilingual{\matriculationyearname}{Immatrikulationsjahr}{Matriculation year:}
\bilingual{\graduationname}{zur Erlangung des akademischen Grades}{to achieve the academic degree}
\bilingual{\refereename}{Gutachter}{Referee}
\bilingual{\advisorname}{Fachreferent}{Advisor}
\bilingual{\supervisorname}{Betreuer}{Supervisor}
\bilingual{\professorname}{Betreuender Hochschullehrer}{Supervising professor}
\bilingual{\datename}{Eingereicht am:}{Submitted on:}
%% Thesis Typen (inspiriert von dem Paket tudscr):
\bilingual{\dissertationname}{Dissertation}{Dissertation} % diss, doctoral, phd 
\bilingual{\diplomathesisname}{Diplomarbeit}{Diploma Thesis} % diploma 
\bilingual{\masterthesisname}{Master-Arbeit}{Master Thesis} % master 
\bilingual{\bachelorthesisname}{Bachelor-Arbeit}{Bachelor Thesis} % bachelor 
\bilingual{\studentthesisname}{Studienarbeit}{Student Thesis} % student 
\bilingual{\studentresearchname}{Großer Beleg}{Student Research Project} % evidence 
\bilingual{\projectpapername}{Projektarbeit}{Project Paper} % project
\bilingual{\seminarpapername}{Seminararbeit}{Seminar Paper} % seminar
\bilingual{\termpapername}{Hausarbeit}{Term Paper} % term
\bilingual{\researchname}{Forschungsbericht}{Research Report} % research 
\bilingual{\logname}{Protokoll}{Log} % log 
\bilingual{\reportname}{Bericht}{Report} % report
\bilingual{\internshipname}{Praktikumsbericht}{Internship Report} % internship 
\bilingual{\lecturename}{Vorlesungsmitschrift}{Lecture Notes} % lecture 
\bilingual{\lectureauthorname}{Mitschrift von}{Notes by} % lecture  
\bilingual{\lessonname}{\"Ubungsmitschrift}{Lesson Notes} % lecture  
\bilingual{\practname}{Praktikumsmitschrift}{Pracital Course Notes} % lecture 
%%% Variablen für die Titelseite (inspiriert von dem Paket tudscr):
\renewcommand{\subject}[1]{\checkThesis{#1}{\@subject}}
\newcommand{\faculty}[1]{\gdef\@faculty{#1}}
\newcommand{\department}[1]{\gdef\@department{#1}}
\newcommand{\institute}[1]{\gdef\@institute{#1}}
\newcommand{\chair}[1]{\gdef\@chair{#1}}
\newcommand{\authormore}[1]{\gdef\@authormore{#1}}
\newcommand{\dateofbirth}[1]{\gdef\@dateofbirth{#1}}
\newcommand{\placeofbirth}[1]{\gdef\@placeofbirth{#1}}
\newcommand{\course}[1]{\gdef\@course{#1}}
\newcommand{\discipline}[1]{\gdef\@discipline{#1}}
\newcommand{\matriculationnumber}[1]{\gdef\@matriculationnumber{#1}}
\newcommand{\matriculationyear}[1]{\gdef\@matriculationyear{#1}}
\newcommand{\thesis}[1]{\checkThesis{#1}{\@thesis}}
\newcommand{\graduation}[2][]{\gdef\@graduation{#2}\gdef\@graduationshort{#1}}
\newcommand{\supervisor}[1]{\gdef\@supervisor{#1}}
\newcommand{\referee}[1]{\gdef\@referee{#1}}
\newcommand{\advisor}[1]{\gdef\@advisor{#1}}
\newcommand{\professor}[1]{\gdef\@professor{#1}}
%%% Titelseite (inspiriert von dem Paket tudscr):
\renewcommand*{\maketitle}{%
\begin{titlepage}%
\begin{flushleft}
\renewcommand{\thefootnote}{$\star$}%
%%% Kopfzeile HTW-Style:
\@htwheader\\[13mm]
%%% Die Zeilen für titlehead, subject, title, subtitle und author sind immer vorhanden (spacing)
{\noindent\@titlehead{}\strut}\\[3mm]%
{\Large\@subject{}\strut}\\[8mm]%
{\Huge\bfseries%
\ifdefined\@nouppercase\@title{}\else\MakeUppercase{\@title{}}\fi%
\strut}\\
{\Large\bfseries%
\ifdefined\@nouppercase\@subtitle{}\else\MakeUppercase{\@subtitle{}}\fi%
\strut}\\[11mm]
\ifdefined\@islecture\strut\else{\Large\@author{}\strut}\fi % Autor wird bei Vorlesungsmitschrift unten angezeigt
%%% Weitere Felder werden nur angezeigt, wenn sie definiert wurden
\ifdefined\@authormore\ifdefined\@islecture\else\\\fi{\@authormore{}\strut}\fi
\ifdefined\@dateofbirth\\\dateofbirthname{} \@dateofbirth{} \ifdefined\@placeofbirth\placeofbirthname{} \@placeofbirth{}\strut\fi\fi
\ifdefined\@course\\\coursename{} \@course{}\strut\fi
\ifdefined\@discipline\\\disciplinename{} \@discipline{}\strut\fi
\ifdefined\@matriculationnumber\\\matriculationnumbername{} \@matriculationnumber{}\strut\fi
\ifdefined\@matriculationyear\\\matriculationyearname{} \@matriculationyear{}\strut\fi
\ifdefined\@thesis\\[10mm]{\LARGE\bfseries
\ifdefined\@nouppercase\@thesis{}\else\MakeUppercase{\@thesis{}}\fi
\strut}
\ifdefined\@graduation\\[2mm]{\graduationname{}\strut}
\\[5mm]{\LARGE\bfseries
\ifdefined\@nouppercase\relax\else\MakeUppercase{\fi%
\@graduation{} \ifdefined\@graduationshort(\@graduationshort{})\fi
\ifdefined\@nouppercase\relax\else}\fi\strut}\fi\fi
\vfill\noindent
\ifdefined\@islecture\\[1mm]{\lectureauthorname{}\strut}\\[.7mm]{\Large\@author{}\strut}\fi
\ifdefined\@referee{\refereename{}\strut}\\[.7mm]{\Large\@referee{}\strut}\fi
\ifdefined\@advisor\\[1mm]{\advisorname{}\strut}\\[.7mm]{\Large\@advisor{}\strut}\fi
\ifdefined\@supervisor\\[1mm]{\supervisorname{}\strut}\\[.7mm]{\Large\@supervisor{}\strut}\fi
\ifdefined\@professor\\[1mm]{\professorname{}\strut}\\[.7mm]{\Large\@professor{}\strut}\fi
\ifdefempty{\@date}{}{\ifdefined\@simpledate\\[6mm]\relax\else\\[3mm]\datename{}~\fi\@date{}\strut}
\ifdefempty{\@publishers}{}{\\[4.5mm]{\Large\@publishers{}\strut}}%
\renewcommand\footnoterule{}
\@thanks\global\let\@thanks\@empty
\end{flushleft}
\end{titlepage}
\setcounter{footnote}{0}
}
\else%
\PackageError{htwcd}{htwcd needs a KOMA class!}%
\fi%
\makeatother
%%% === Ende Paket htwscr ===
%-- Nummerierung bis Subsubsection für Report
\setcounter{secnumdepth}{3}	% zählt auch subsubsection
\setcounter{tocdepth}{3}	% Inhaltsverzeichnis bis in subsubsection

% LuaLaTeX Unterscheidung
\usepackage{ifluatex}

% Input inkl. Umlaute, Silbentrennung
\ifluatex
	\usepackage{fontspec}
\else
	\usepackage[utf8]{inputenc}
	\usepackage[T1]{fontenc}
\fi
\usepackage[english,ngerman]{babel}
\usepackage{csquotes}	% Anführungszeichen
\usepackage{marvosym}
\usepackage{eurosym}

% Style-Aufhübschung
\usepackage{soul, color}	% Kapitälchen, Unterstrichen, Durchgestrichen usw. im Text
%\usepackage{titleref}
\usepackage[breakwords, fit]{truncate}	% Abschneiden von Sätzen
\renewcommand{\TruncateMarker}{\,…}

% Mathe usw.
\usepackage{amssymb}
\usepackage{amsthm}
\usepackage[fleqn,intlimits]{amsmath}	% fleqn: align-Umgebung rechtsbündig; intlimits: Integralgrenzen immer ober-/unterhalb
%\usepackage{mathtools} % u.a. schönere underbraces
\usepackage{xcolor}
\usepackage{esint}	% Schönere Integrale, \oiint vorhanden
\everymath=\expandafter{\the\everymath\displaystyle}	% Mathe Inhalte werden weniger verkleinert
\usepackage{wasysym}	% mehr Symbole, bspw \lightning
%\renewcommand{\int}{\int\limits}
%\usepackage{xfrac}	% mehr fracs: sfrac{}{}
\let\oldemptyset\emptyset	% schöneres emptyset
\let\emptyset\varnothing
%\RequirePackage{mathabx}	% mehr Symbole
\mathchardef\mhyphen="2D	% Hyphen in Math
\usepackage{mathtools}		% zum Beispiel {pmatrix}[r] für Formatierung der Matrix

% tikz usw.
\usepackage{tikz}
\usepackage{pgfplots}
\pgfplotsset{compat=1.11}	% Umgeht Fehlermeldung
\usetikzlibrary{graphs}
%\usetikzlibrary{through}	% ???
\usetikzlibrary{arrows}
\usetikzlibrary{arrows.meta}	% Pfeile verändern / vergrößern: \draw[-{>[scale=1.5]}] (-3,5) -> (-3,3);
\usetikzlibrary{automata,positioning} % Zeilenumbruch im Node node[align=center] {Text\\nächste Zeile} automata für Graphen
\usetikzlibrary{matrix}
\usetikzlibrary{patterns}	% Schraffierte Füllung
\usetikzlibrary{shapes.geometric}	% Polygon usw.
\tikzstyle{reverseclip}=[insert path={	% Inverser Clip \clip
	(current page.north east) --
	(current page.south east) --
	(current page.south west) --
	(current page.north west) --
	(current page.north east)}
% Nutzen: 
%\begin{tikzpicture}[remember picture]
%\begin{scope}
%\begin{pgfinterruptboundingbox}
%\draw [clip] DIE FLÄCHE, IN DER OBJEKT NICHT ERSCHEINEN SOLL [reverseclip];
%\end{pgfinterruptboundingbox}
%\draw DAS OBJEKT;
%\end{scope}
%\end{tikzpicture}
]	% Achtung: dafür muss doppelt kompliert werden!
\usepackage{graphpap}	% Grid für Graphen
\tikzset{every state/.style={inner sep=2pt, minimum size=2em}}
\usetikzlibrary{mindmap, backgrounds}
%\usepackage{tikz-uml}	% braucht Dateien: http://perso.ensta-paristech.fr/~kielbasi/tikzuml/

% Tabular
\usepackage{longtable}	% Große Tabellen über mehrere Seiten
\usepackage{multirow}	% Multirow/-column: \multirow{2[Anzahl der Zeilen]}{*[Format]}{Test[Inhalt]} oder \multicolumn{7[Anzahl der Reihen]}{|c|[Format]}{Test2[Inhalt]}
\renewcommand{\arraystretch}{1.3} % Tabellenlinien nicht zu dicht
\usepackage{colortbl}
\arrayrulecolor{gray}	% heller Tabellenlinien
\usepackage{array}	% für folgende 3 Zeilen (für Spalten fester breite mit entsprechender Ausrichtung):
\newcolumntype{L}[1]{>{\raggedright\let\newline\\\arraybackslash\hspace{0pt}}m{\dimexpr#1\columnwidth-2\tabcolsep-1.5\arrayrulewidth}}
\newcolumntype{C}[1]{>{\centering\let\newline\\\arraybackslash\hspace{0pt}}m{\dimexpr#1\columnwidth-2\tabcolsep-1.5\arrayrulewidth}}
\newcolumntype{R}[1]{>{\raggedleft\let\newline\\\arraybackslash\hspace{0pt}}m{\dimexpr#1\columnwidth-2\tabcolsep-1.5\arrayrulewidth}}
\usepackage{booktabs}	% Kompatibilität für array und \hrule
\usepackage{caption}	% Um auch unbeschriftete Captions mit \caption* zu machen

% Nützliches
\usepackage{verbatim}	% u.a. zum auskommentieren via \begin{comment} \end{comment}
\usepackage{tabto}	% Tabs: /tab zum nächsten Tab oder /tabto{.5 \CurrentLineWidth} zur Stelle in der Linie
\NumTabs{6}	% Anzahl von Tabs pro Zeile zum springen
\usepackage{listings} % Source-Code mit Tabs
\usepackage{lstautogobble} 
\usepackage{enumitem}	% Anpassung der enumerates
%\setlist[enumerate,1]{label=(\arabic*)}	% global andere Enum-Items
\renewcommand{\labelitemiii}{$\scriptscriptstyle ^\blacklozenge$} % global andere 3. Item-Aufzählungszeichen
\usepackage{letltxmacro} % neue Definiton von Grundbefehlen
% Nutzen:
%\LetLtxMacro{\oldemph}{\emph}
%\renewcommand{\emph}[1]{\oldemph{#1}}
\RequirePackage{xpatch}	% ua. Konkatenieren von Strings/Variablen (etoolbox)
\usepackage{xstring}	% String Operationen
\usepackage{minibox}	% Minibox anstatt \fbox{} für Boxen mit (manuellen) Zeilenumbruch
\usepackage{framed}		% \begin{framed} statt minibox/fbox mit automatischem Zeilenumbruch
\usepackage{microtype}	% Verfeinerte Hyphenation
\usepackage{import}		
\usepackage{standalone}	% Zusammen mit import: Einfügen von separaten Latex-Dateien (die schon header haben)

% Einrichtung von lst
%\renewcommand*{\ttdefault}{qcr}
\lstset{
basicstyle=\ttfamily, 
%basicstyle=\footnotesize\sffamily\color{black},
%mathescape=true, 
%escapeinside=^^, 
autogobble, 
tabsize=2,
frame=single,
rulecolor=\color{lightgray},
numbers=left,
numbersep=5pt,
numberstyle=\tiny\color{gray},
commentstyle=\color{gray},
keywordstyle=\color{htwblue},
stringstyle=\color{htworange},
morecomment=[l][\color{magenta}]{\#}
showspaces=false,
showstringspaces=false,
breaklines=true,
literate=%
    {Ö}{{\"O}}1
    {Ä}{{\"A}}1
    {Ü}{{\"U}}1
    {ß}{{\ss}}1
    {ü}{{\"u}}1
    {ä}{{\"a}}1
    {ö}{{\"o}}1
    {„}{{\glqq}}1
    {“}{{\grqq}}1
    {~}{{\textasciitilde}}1
}
\usepackage{scrhack} % Fehler umgehen
\def\ContinueLineNumber{\lstset{firstnumber=last}} % vor lstlisting. Zum wechsel zum nicht-kontinuierlichen muss wieder \StartLineAt1 eingegeben werden
\def\StartLineAt#1{\lstset{firstnumber=#1}} % vor lstlisting \StartLineAt30 eingeben, um bei Zeile 30 zu starten
\let\numberLineAt\StartLineAt

% BibTeX
\usepackage[
backend=bibtex,
bibencoding=ascii,
style=authoryear, citestyle=authoryear-ibid,
%doi=false,
%isbn=false,
%url=false
]{biblatex}	% BibTeX
\usepackage{makeidx}
%\makeglossary
%\makeindex
\usepackage{url}	% Darstellung von URLs

% Grafiken
\usepackage{graphicx}
\usepackage{epstopdf}	% eps-Vektorgrafiken einfügen
\usepackage{transparent}	% transparent nutzen: {\transparent{0.4} ...}
%\epstopdfsetup{outdir=\customDir}
% Prüft, ob Grafik existiert (mit \ifvalidimage{}{}) [Quelle: https://tex.stackexchange.com/a/99176]:
\makeatletter
\newif\ifgraphicexist
\catcode`\*=11
\newcommand\ifvalidimage[1]{%
    \begingroup
    \global\graphicexisttrue
    \let\input@path\Ginput@path
    \filename@parse{#1}%
    \ifx\filename@ext\relax
    \@for\Gin@temp:=\Gin@extensions\do{%
        \ifx\Gin@ext\relax
        \Gin@getbase\Gin@temp
        \fi}%
    \else
    \Gin@getbase{\Gin@sepdefault\filename@ext}%
    \ifx\Gin@ext\relax
    \global\graphicexistfalse
    \def\Gin@base{\filename@area\filename@base}%
    \edef\Gin@ext{\Gin@sepdefault\filename@ext}%
    \fi
    \fi
    \ifx\Gin@ext\relax
    \global\graphicexistfalse
    \else 
    \@ifundefined{Gin@rule@\Gin@ext}%
    {\global\graphicexistfalse}%
    {}%
    \fi  
    \ifx\Gin@ext\relax 
    \gdef\imageextension{unknown}%
    \else
    \xdef\imageextension{\Gin@ext}%
    \fi 
    \endgroup 
    \ifgraphicexist
    \expandafter \@firstoftwo
    \else
    \expandafter \@secondoftwo
    \fi 
} 
\catcode`\*=12
\makeatother
\LetLtxMacro{\forceincludegraphics}{\includegraphics}	% neuer Befehl für includegraphics
\renewcommand{\includegraphics}[2][]{	% altes includegraphics neu definieren, damit es auch nicht vorhandene einfügt
\ifvalidimage{#2}{
\forceincludegraphics[#1]{#2}
}{
\message{Achtung: Grafik wurde nicht gefunden: '#2'}
\minibox[frame]{
\textbf{\StrSubstitute{#2}{_}{\_}}%  \ifthenelse{\isempty{#1}}{}{\\\textit{#1}}
}}}

% pdf-Setup
\usepackage{pdfpages}
\usepackage[bookmarks,%
bookmarksopen=false,% Klappt die Bookmarks in Acrobat aus
colorlinks=true,%
linkcolor=black,%
citecolor=black,%
urlcolor=htworange,%
]{hyperref}

% Nach Dokumentenbeginn ausführen:
\makeatletter
\AtBeginDocument{	
	% pdf Informationen
	\hypersetup{
		pdfauthor={\@author},
		pdftitle={\@title},
		pdfproducer={LaTeX Layout erstellt von Falk-Jonatan Strube}
	}
	% signierte Fußzeile
	%\providecommand{\ofootname}{Mitschrift von}
	\providecommand{\pathtomaster}{../}
	\ofoot{%
	%\ifdefined\ofoottext
	%\else
	%\begin{minipage}{40mm}
	%\begin{flushright}
	%\textcolor{gray}{
	%\ofootname\\
	%Falk-Jonatan Strube\vspace*{-1mm}}
	%\end{flushright}
	%\end{minipage}
	%\fi%
	%\hspace*{1mm}%
	\begin{minipage}{10 mm}
	\includegraphics[height=9mm]{\pathtomaster/FJS-logo_gray}
	\end{minipage}%
	}
}
\makeatother



%%% ==================================================
%%% ===================== MACROS =====================
%%% ==================================================
% Horizontale Linie:
\newcommand{\HRule}[1][\medskipamount]{\par
  \vspace*{\dimexpr-\parskip-\baselineskip+#1}
  \noindent\rule[0.2ex]{\linewidth}{0.2mm}\par
  \vspace*{\dimexpr-\parskip-.5\baselineskip+#1}}
% Gestrichelte horizontale Linie:
\RequirePackage{dashrule}
\newcommand{\HDRule}[1][\medskipamount]{\par
  \vspace*{\dimexpr-\parskip-\baselineskip+#1}
  \noindent\hdashrule[0.2ex]{\linewidth}{0.2mm}{1mm} \par
  \vspace*{\dimexpr-\parskip-.5\baselineskip+#1}}
% Mathe in Anführungszeichen:
\newsavebox{\mathbox}\newsavebox{\mathquote}
\makeatletter
\newcommand{\mq}[1]{% \mathquotes{<stuff>}
  \savebox{\mathquote}{\text{"}}% Save quotes
  \savebox{\mathbox}{$\displaystyle #1$}% Save <stuff>
  \raisebox{\dimexpr\ht\mathbox-\ht\mathquote\relax}{"}#1\raisebox{\dimexpr\ht\mathbox-\ht\mathquote\relax}{''}
}
\makeatother

% Paragraph mit Zähler (Section-Weise)
\newcounter{cparagraphC}
\newcommand{\cparagraph}[1]{
\stepcounter{cparagraphC}
\paragraph{\thesection{}-\thecparagraphC{} #1}
%\addcontentsline{toc}{subsubsection}{\thesection{}-\thecparagraphC{} #1}
\label{\thesection-\thecparagraphC}
}
\makeatletter
\@addtoreset{cparagraphC}{section}
\makeatother


% (Vorlesungs-)Folien einbinden:
\makeatletter
% Folien-Pfad (inkl. Dateiname ohne Endung und ggf. ohne Nummerierung)
\newcommand{\@slidePath}{}
\newcommand{\slidePath}[1]{\renewcommand{\@slidePath}{#1}}
% Folien Eigenschaften
\newcommand{\@slideScale}{0.5}
\newcommand{\slideScale}[1]{\renewcommand{\@slideScale}{#1}}
\AtBeginDocument{
% Folien von einer Datei skaliert
\newcommand{\slide}[2][\@slideScale]{\slides[#1]{}{#2}}
\newcommand{\slideTrim}[6][\@slideScale]{\slides[#1 , clip,  trim = #5cm #4cm #6cm #3cm]{}{#2}}
% Folien von mehreren nummerierten Dateien skaliert
\newcommand{\slides}[3][\@slideScale]{\begin{center}
\includegraphics[page=#3, scale=#1]{\@slidePath #2.pdf}
\end{center}}
}
\makeatother

% \emph{} anders definieren
\makeatletter
\DeclareRobustCommand{\em}{%
  \@nomath\em \if b\expandafter\@car\f@series\@nil
  \normalfont \else \scshape \fi}
\makeatother

% unwichtiges
\newcommand{\unimptnt}[1]{{\transparent{0.5}#1}}

% alph. enumerate
\newenvironment{anumerate}{\begin{enumerate}[label=(\alph*)]}{\end{enumerate}} % Alphabetische Aufzählung

% Hanging parameters
\newcommand{\hangpara}[1]{\par\noindent\hangindent+2em\hangafter=1 #1\par\noindent}

%% EINFACHE BEFEHLE

% Abkürzungen Mathe
\newcommand{\EE}{\mathbb{E}}
\newcommand{\QQ}{\mathbb{Q}}
\newcommand{\RR}{\mathbb{R}}
\newcommand{\CC}{\mathbb{C}}
\newcommand{\NN}{\mathbb{N}}
\newcommand{\ZZ}{\mathbb{Z}}
\newcommand{\PP}{\mathbb{P}}
\renewcommand{\SS}{\mathbb{S}}
\newcommand{\cA}{\mathcal{A}}
\newcommand{\cB}{\mathcal{B}}
\newcommand{\cC}{\mathcal{C}}
\newcommand{\cD}{\mathcal{D}}
\newcommand{\cE}{\mathcal{E}}
\newcommand{\cF}{\mathcal{F}}
\newcommand{\cG}{\mathcal{G}}
\newcommand{\cH}{\mathcal{H}}
\newcommand{\cI}{\mathcal{I}}
\newcommand{\cJ}{\mathcal{J}}
\newcommand{\cM}{\mathcal{M}}
\newcommand{\cN}{\mathcal{N}}
\newcommand{\cP}{\mathcal{P}}
\newcommand{\cR}{\mathcal{R}}
\newcommand{\cS}{\mathcal{S}}
\newcommand{\cZ}{\mathcal{Z}}
\newcommand{\cL}{\mathcal{L}}
\newcommand{\cT}{\mathcal{T}}
\newcommand{\cU}{\mathcal{U}}
\newcommand{\cX}{\mathcal{X}}
\newcommand{\cV}{\mathcal{V}}
\renewcommand{\phi}{\varphi}
\renewcommand{\epsilon}{\varepsilon}
\renewcommand{\theta}{\vartheta}

% Verschiedene als Mathe-Operatoren
\DeclareMathOperator{\arccot}{arccot}
\DeclareMathOperator{\arccosh}{arccosh}
\DeclareMathOperator{\arcsinh}{arcsinh}
\DeclareMathOperator{\arctanh}{arctanh}
\DeclareMathOperator{\arccoth}{arccoth} 
\DeclareMathOperator{\var}{Var} % Varianz 
\DeclareMathOperator{\cov}{Cov} % Co-Varianz 

% Farbdefinitionen
\definecolor{red}{RGB}{180,0,0}
\definecolor{green}{RGB}{75,160,0}
\definecolor{blue}{RGB}{0,75,200}
\definecolor{orange}{RGB}{255,128,0}
\definecolor{yellow}{RGB}{255,245,0}
\definecolor{purple}{RGB}{75,0,160}
\definecolor{cyan}{RGB}{0,160,160}
\definecolor{brown}{RGB}{120,60,10}

\definecolor{itteny}{RGB}{244,229,0}
\definecolor{ittenyo}{RGB}{253,198,11}
\definecolor{itteno}{RGB}{241,142,28}
\definecolor{ittenor}{RGB}{234,98,31}
\definecolor{ittenr}{RGB}{227,35,34}
\definecolor{ittenrp}{RGB}{196,3,125}
\definecolor{ittenp}{RGB}{109,57,139}
\definecolor{ittenpb}{RGB}{68,78,153}
\definecolor{ittenb}{RGB}{42,113,176}
\definecolor{ittenbg}{RGB}{6,150,187}
\definecolor{itteng}{RGB}{0,142,91}
\definecolor{ittengy}{RGB}{140,187,38}

% Textfarbe ändern
\newcommand{\tred}[1]{\textcolor{red}{#1}}
\newcommand{\tgreen}[1]{\textcolor{green}{#1}}
\newcommand{\tblue}[1]{\textcolor{blue}{#1}}
\newcommand{\torange}[1]{\textcolor{orange}{#1}}
\newcommand{\tyellow}[1]{\textcolor{yellow}{#1}}
\newcommand{\tpurple}[1]{\textcolor{purple}{#1}}
\newcommand{\tcyan}[1]{\textcolor{cyan}{#1}}
\newcommand{\tbrown}[1]{\textcolor{brown}{#1}}

% Umstellen der Tabellen Definition
\newcommand{\mpb}[1][.3]{\begin{minipage}{#1\textwidth}\vspace*{3pt}}
\newcommand{\mpe}{\vspace*{3pt}\end{minipage}}

\newcommand{\resultul}[1]{\underline{\underline{#1}}}
\newcommand{\parskp}{$ $\\}	% new line after paragraph
\newcommand{\parnl}{\unskip\parfillskip 0pt \par} % justified new line (brace line!)

% mathematische Abkürzungen
\newcommand{\corr}{\;\widehat{=}\;}
\newcommand{\mdeg}{^{\circ}}
\newcommand{\nok}[2]{\binom{#1}{#2}}	% n über k BESSER: \binom{n}{k}
\newcommand{\mtr}[1]{\begin{pmatrix}#1\end{pmatrix}}	% Matrix
\newcommand{\dtr}[1]{\begin{vmatrix}#1\end{vmatrix}}	% Determinante (Betragsmatrix)
\LetLtxMacro{\originalVec}{\vec}
\renewcommand{\vec}[1]{\underline{#1}}	% Vektorschreibweise
\newcommand{\imptnt}[1]{\colorbox{red!30}{#1}}	% Wichtiges
\newcommand{\intd}[1]{\,\mathrm{d}#1}
\newcommand{\diffd}[1]{\mathrm{d}#1}
% für Module-Rechnung: \pmod{}
\newcommand{\unit}[1]{\,\mathrm{#1}}
%\LetLtxMacro{\ntilde}{\tilde}
%\renewcommand{\tilde}{\widetilde}

% weiter Abkürzungen
\newcommand{\gdw}{genau dann wenn}
\newcommand{\lecdate}[1]{\marginpar{\small\textcolor{gray}{Vorlesung\newline #1}}}

\thesis{Entwicklerdokumentation -- Grobentwurf}
\date{}

\makeatletter
\gdef\@islecture{\relax}
\makeatother

\faculty{Fakultät Informatik/Mathematik}
\chair{Lehrstuhl für Softwaretechnologie}
\subject{Software Engineering 2}
\title{LED-Leuchtturm}
%\subtitle{Entwicklung eines LED-Leuchtturms zur Anzeige des aktuellen Maschinenzustands}
\bilingual{\lectureauthorname}{Gruppenmitglieder}{Group members}
\author{
\texorpdfstring{Ragnar Luga		\tabto{.35\linewidth} (s75594)\\
Raphael Pour									\tabto{.35\linewidth} (s74020)\\
Torsten Polczyk								\tabto{.35\linewidth} (s74097)\\
Simon Retsch									\tabto{.35\linewidth} (s74077)\\
Eric Schmidtgen								\tabto{.35\linewidth} (s73952)\\
Ben Schönherr								\tabto{.35\linewidth} (s74083)\\
Falk-Jonatan Strube						\tabto{.35\linewidth} (s74053)\\
Jürgen Tomaszewski						\tabto{.35\linewidth} (s76451)\\\vspace*{-1em}}{
Ragnar Luga,
Raphael Pour,
Torsten Polczyk,
Simon Retsch,
Erik Schmidtgen,
Ben Schoenherr,
Falk-Jonatan Strube,
Jürgen Tomaszewski
}
}
\bilingual{\supervisorname}{Auftraggeber}{Client}
\supervisor{Dipl.-Inf. Robert Ringel}
\professor{Prof. Dr. Anna Sabine Hauptmann}

\let\tableofcontentsORIG\tableofcontents
\renewcommand\tableofcontents{\tableofcontentsORIG\clearpage}

\begin{document}
\maketitle
\tableofcontents

Die im Grobentwurf erarbeiteten Themen geben die Eckpunkte für das Softwaresystem vor. Das Pflichtenheft wurde auf dessen Basis formuliert.

\subsection{Zielsetzung}
Die Zielsetzung für das Softwaresystem \texttt{ledleuchtturm} ergab sich aus der Aufgabenstellung. Diese war Grundlage der Entwicklung und wurde vom Kunden (Herrn Dipl.-Inf. Robert Ringel) bestätigt und erläutert:
\begin{quote}
„Entwicklung eines LED‐Leuchtturms zur Anzeige des aktuellen Maschinenzustands:\\
Mit Hilfe eines Raspberry-Pi-Computers ist der aktuelle Maschinenzustand einer Maschine der Cinderella-Modellfabrik in Form einer LED-Leuchtturmanzeige zu visualisieren. Der Maschinenzustand ist aus der
PostgreSQL-Datenbank des Produktionssteuerungssystems zu entnehmen.“
\begin{flushright}
\textit{Belegarbeit „Software-Engineering II“: Sommersemester 2017}
\end{flushright}
\end{quote}

\subsection{Systemeinsatz und Umgebung}
Im Kundengespräch wurde die Festlegung, dass das Softwaresystem auf einem Raspberry Pi der zweiten Generation laufen soll, vereinbart. Dieser Computer soll von einer Datenbank, die sich im selben Netzwerk befindet, Informationen über den aktuellen Status einer Maschine erhalten. Das Übertragen der Daten ist in der Datenbank von einem Trigger auszulösen. Die LEDs für die Anzeige des Status sind direkt an den GPIO-Port anzuschließen.

\subsection{Benutzerschnittstellen}
Im Hinblick darauf, dass das Softwaresystem stets im Hintergrund läuft, beschränkt sich die Benutzerschnittstelle auf das Starten und Beenden der Software. Beim Starten ist die die Konfiguration der Software relevant, die jedoch in eine Datei ausgelagert und nicht durch ein Nutzerinterface eingegeben werden soll.

\subsection{Funktionale Anforderungen}
Funktionale Anforderungen sind teilweise direkt durch die Aufgabenstellung bedingt: Die Software soll einen Status(-wechsel) der Datenbank durch LEDs anzeigen. Die Analyse der Benutzerschnittstelle hat ergeben, dass das Softwaresystem weiterhin eine Konfigurationsdatei besitzen und diese auswerten muss. In der Konfigurationsdatei müssen alle notwendigen Informationen stehen, um sich mit der Datenbank zu verbinden, den Status der gewünschten Maschine auszulesen und auf dem im Trigger eingerichteten Port auf einen Statuswechsel der Datenbank zu warten.

\subsection{Qualitätsanforderungen}
Die Qualitätsanforderungen des Softwaresystems beschränken sich auf zwei Punkte:
\begin{enumerate}
	\item Das Softwaresystem hat innerhalb einer Sekunde auf einen Statuswechsel in der Datenbank zu reagieren.
	\item Das Softwaresystem soll nach dem Starten selbstständig laufen.
\end{enumerate}
Dabei wird mit der ersten Anforderung betont, dass das System annäherungsweise in Echtzeit laufen soll. Da die Hardware allerdings keine Laufzeit in Echtzeit hergeben kann, wurde die eine Sekunde als Schwellwert definiert.

\subsection{Rahmenbedingungen}
Das Softwaresystem kann nur lauffähig sein, wenn der entsprechende Trigger und die Konfigurationsdatei vom Nutzer selbständig korrekt eingerichtet wurden. Die Benutzerdokumentation hat das Einrichten der Software ausreichend zu beschreiben.

\end{document}
% Grobentwurf aus Pflichtenheft entnehmen und noch mal als Dokumentation aufarbeiten.

\clearpage
\section{Feinentwurf}
\label{feinentw}
%\subsection{Einleitung}
	In diesem Abschnitt wird der konzeptionelle Feinentwurf beschrieben. Er beinhaltet alle Diagramme, Richtlinien, Hinweise und Dateien, die notwendig für die Implementierung sind. Aufgabe des Entwurfes ist es, die Implementation und Entwicklung der geforderten Software mit ihren Rahmenbedingungen und Anforderungen in ihrer hinreichenden Komplexität handhabbar zu gestalten. Somit wird auf alle wichtigen Punkte, Sachverhalte und Probleme eingegangen, die in Verbindung zum Erstellen des Produktes stehen.
	
	Als Grundlage für den Entwurf wird das Pflichtenheft genommen, welches von Eric Schmidtgen (Analyst) erstellt wurde und von Jürgen Tomaszewski (Projektleiter) sowie unserem Auftraggeber unterschrieben wurde. 
\subsection{Programmierstandards}	
\subsubsection{Programmiersprache}	
	Da nicht nur der Implementierer, sondern mindestens auch der Qualitätssichernde, der Datenbankbeauftragte und der Projektleiter den Quellcode verstehen muss, ist es wichtig, dass sich gemeinsam auf eine Programmiersprache geeinigt wird. Zur Auswahl standen die Programmiersprachen C, Java und Python. Das Programm \texttt{ledleuchtturm} wird mit der Programmiersprache Python geschrieben, da mit dieser Programmiersprache Dateiarbeit und die Anbindung zur Datenbank relativ leicht umsetzbar ist. Die Version 2.7.0. ist auf den meisten Linuxsystemen standardisiert. 
\subsubsection{Quellcode}
	Der Programmcode ist sauber und verständlich zu halten. Da die Qualität des Quellcodes von Raphael Pour (Qualitätsmanager) überprüft wird, sind hier nur erste Richtlinien für den Quellcode genannt:
	
	Variablennamen und Kommentare sind im Quellcode in Englisch zu verfassen. Jede Funktion im Quellcode \emph{muss} einen Rückgabewert besitzen, damit geprüft werden kann, ob die Funktionen korrekt abgearbeitet wurden. Erst nachdem man den Rückgabewert geprüft hat, werden die Ausgaben (siehe \autoref{entw:ausg}) gemacht.
	
	Der Implementierer darf im Quellcode Funktionen ändern oder hinzufügen, soweit es die funktionalen Anforderungen (siehe \autoref{entw:funktAnf}) nicht beeinträchtigt.
\clearpage
\subsection{Diagramme}
ERM, Klassendiagramm und Anwendungsfalldiagramm entfallen, da sie in diesem Kontext nicht gewinnbringend sind: 
\begin{itemize}
\item Das Programm ist aufgrund des relativ geringen Umfangs nicht objektorientiert. Es wird also kein Klassendiagramm benötigt. Dementsprechend entfällt auch das ERM.
\item Das Programm beschränkt sich auf einen Anwendungsfall (Programm wird gestartet).
\end{itemize}
Somit besteht der Entwurf ausschließlich aus folgenden Elementen:
\begin{itemize}
	\item Aktivitätsdiagramm des vollständigen Programms, wie es im Pflichtenheft festgehalten wurde (\autoref{entw:aktdia}) mit funktionalen Anforderungen (\autoref{entw:funktAnf}) und Ausgaben (\autoref{entw:ausg})
	\item Schaltskizze (\autoref{entw:schltsk})
	\item Zustandsdiagramm (\autoref{entw:zstdia})
	\item Konfigurationsdatei mit einem Beispiel (\autoref{entw:konf})
\end{itemize}	

\clearpage
\subsubsection{Aktivitätsdiagramm}
	Das Aktivitätsdiagramm beschreibt den regulären Ablauf des Programms vom Start bis zum Beenden. Grundsätzlich bestand hier die Herausforderung darin alle notwendigen funktionalen Anforderungen in das Aktivitätsdiagramm einzubringen und gleichzeitig nur den Programmablauf zu beachten. %Es werden keine mögliche Implementationslösungen betrachtet.
\label{entw:aktdia}
\begin{center}
\includegraphics[width=.99\textwidth , trim={0.5cm 2.1cm 0.4cm 1cm},clip]{include/Aktivitaetsdiagramm_include}
\end{center}

\clearpage
\subsubsection{Funktionale Anforderungen}
\label{entw:funktAnf}
	Die unten genannten funktionalen Anforderungen \emph{müssen} in unserer Software umgesetzt werden.
\begin{center}
\resizebox{.99\textwidth}{!}{
\begin{tabular}{ | C{.13} | L{.35} | L{.2} | L{.2}  | L{.17} | }
	\hline
	Kennung	& Funktion	& Eingangsdaten	& Ausgangsdaten & Ausgabe \\ \hline\hline
	$F_{1}$ & Prüfe, ob eine Konfigurations- datei mit Pfad angehängt ist	& Angehängte Konfigurationsdatei	& Konfigurationsdatei	&	 $A_{1},A_{2}$\\ \hline
	
	$F_{2}$ & Prüfe, ob die Konfigurations- datei gefunden wird 	& Konfigurationsdatei	& & $A_{3},A_{4.1},A_{4.2}$\\ \hline
	$F_{3}$ & Prüfe, ob es sich um die Konfigurationsdatei handelt	& Konfigurationsdatei	& 	& $A_{5}$,$A_{6}$\\ \hline
	$F_{4}$ & Lade Datenbankinformationen aus der Konfigurationsdatei 	& Konfigurationsdatei& Datenbankinformationen & \\ \hline
	$F_{5}$ & Teste die Verbindung zur Datenbank & Datenbankinformationen	&	& $A_{7},A_{8}$\\ \hline
	$F_{6}$ & Prüfe, ob Nachricht von Datenbank-Trigger eingetroffen & Datenbank-Trigger	&		&\\ \hline
	$F_{7}$ & Lade den aktuellen Status der	Maschine aus der Datenbank & & Status	&\\ \hline
	$F_{8}$ & Vergleiche Maschinenname des Triggers mit dem Maschinennamen der Konfigurationsdatei & Maschinenname &		&\\ \hline
	$F_{9}$ &Ansteuern der LEDs Entsprechend der Vorgabe & Status	& Status	& $A_{9}$\\ \hline
	$F_{10}$ &Prüfe, ob Abbruchbefehl eingegeben wurde& Abbruchbefehl	&	&$A_{10}$\\ \hline
	$F_{11}$ &Erstelle Konfigurationsdatei mit Struktur, falls kein Pfad angehängt wurde& 	&	&\\\hline
\end{tabular}
}
\end{center}

\clearpage
\subsubsection{Ausgaben}
\label{entw:ausg}
	Hier werden die Nachrichten, die das Programm auf der Konsole ausgibt, aufgelistet. Diese Tabelle soll die unterschiedlichen Meldearten darstellen, die der Implementierer dann entsprechend umzusetzen hat. Alle Nachrichten an den Nutzer werden in Deutsch verfasst. 
\begin{center}
	\begin{tabular}{ | C{.2} | L{.5} | L{.25} | }
	\hline
	Kennung & Nachricht	& Art \\ \hline\hline
	$A_{1}$	& Verwende den Standardpfad für Konfigurationsdatei	& Information	\\ \hline
	$A_{2}$	& Verwende angehängten Pfad mit Konfigurationsdatei	& Information	\\ \hline
	$A_{3}$	& Konfigurationsdatei gefunden	& Info	\\ \hline
	$A_{4.1}$	& Konfigurationsdatei nicht gefunden	& Fehler	\\ \hline
	$A_{4.2}$	& Konfigurationsdatei nicht gefunden, erstelle Muster-Konfigurationsdatei	& Fehler	\\ \hline
	$A_{5}$	& Handelt sich nicht um Konfigurationsdatei	& Fehler	\\ \hline
	$A_{6}$	& Verarbeite Konfigurationsdatei	& Info	\\ \hline
	$A_{7}$	& Verbindung zur Datenbank fehlgeschlagen	& Fehler 	\\ \hline
	$A_{8}$	& Verbindung zur Datenbank	Erfolgreich  Start Abbruchbefehl	& Info	\\ \hline
	$A_{9}$	& Statusänderung, Farbe, Zeit	& Status	\\ \hline
	$A_{9}$	& Programm wird beendet	& Status	\\ \hline	
\end{tabular}
\end{center}

\clearpage
\subsubsection{Schaltskizze}
Im folgendem wird eine Skizze betrachtet. Das  heißt der Aufbau muss im Realen nicht dieser Skizze entsprechen. Die LEDs müssen jedoch einen entsprechenden Vorwiderstand haben. Die Pinbelegung am RaspberryPi ist variabel und kann später über die Konfigurationsdatei (\autoref{entw:konf}) angepasst werden.
\label{entw:schltsk}
\begin{center}
\includegraphics[scale=1]{include/Schaltskizze}
\end{center}

\clearpage
\subsubsection{Zustandsdiagramm}
\label{entw:zstdia}
Das Zustandsdiagramm beschreibt alle Zustände, die ausgehend von der Datenbank eintreten können. Die Software hat fünf wechselbare Zustände, die über die drei LEDs sichtbar gemacht werden können. Der Implementierer hat diese Zustände genauso im Programm umzusetzen.
\begin{center}
\includegraphics[width=.99\textwidth , trim={1.9cm 1.7cm 1.9cm 3.5cm},clip]{include/Zustandsdiagramm_include}
\end{center}

\clearpage
\subsubsection{Konfigurationsdatei}
\label{entw:konf}
	In der Konfigurationsdatei werden alle Variablen und Informationen, die zur Lauffähigkeit des Programms notwendig sind, gespeichert. Die Einhaltung von Standards sind hier besonders wichtig, da das Programm die Datei an einem bestimmten Ort öffnet und den Inhalt entsprechend filtern wird.
\subsubsection*{Richtlinien}
\begin{itemize}
	\item Standard-Konfigurationsdatei heißt Std-Conf-LedLeuchtturm.json
	\item Standard-Konfigurationsdatei liegt im selben Verzeichnis wie das Programm
\end{itemize}

\subsubsection*{Beispiel}
\label{entw:konfbsp}
Die Konfigurationsdatei Std-Conf-LedLeuchtturm.json kann von der Software selbst erstellt werden. Siehe $F_11$ (\autoref{entw:funktAnf}). 
\lstinputlisting[%frame=none
, numbers=none]{include/Std-Conf-LedLeuchtturm.json}
% Feinentwurf aus Entwurf/Entwurf.tex noch mal aufarbeiten und Dokumentieren.

\clearpage
\section{Struktur des Codes}
\label{strukt}
% Struktur des Codes (wie er abgegeben wurde) festhalten.
Das Softwaresystem \texttt{ledleuchtturm} wurde in Python realisiert und ist rein prozedural programmiert. Der Programmcode ist auf fünf Dateien verteilt. Bei vier dieser Dateien handelt es sich um Module mit Funktionen, die zu einem speziellen Aufgabenbereich gehören. Die fünfte Datei ist die Haupt-Datei, die auch beim Ausführen der Software aufgerufen werden muss.
% Wie sind die Dateien/Modul(e) organisiert?

\subsection{Überblick}
% Am besten (statt Klassendiagramm) ein "Datei-Diagramm" erstellen.
Für einen Überblick über die Struktur der Dateien des Softwaresystems mit ihren Funktionalitäten und Variablen sei folgendes Diagramm gegeben:
\begin{center}
	\includegraphics[width=1.0\linewidth]{include/Entwicklerdokumentation_aufbau}
\end{center}

\subsection{Beschreibung}
Aus dem Diagramm lassen sich bereits sehr gut die Funktionalitäten der Software ablesen. Für einen groben Überblick sei der Zweck der Dateien im Folgenden grob umschrieben. Eine detaillierte Beschreibung der Funktionen und Variablen ist im \nameref{epydoc} (\autoref{epydoc}) zu finden.

% Alle Dateien beschreiben: Überblick verschaffen.
\subsubsection*{ledleuchtturm.py}
Die Datei ledleuchtturm.py ist die Haupt-Datei des Projektes. In ihr werden alle Module importiert und die notwendigen Funktionen ausgeführt. Der Code beschränkt sich im Wesentlichen auf das Ausführen von Funktionen, die Auswertung der Rückgabewerte und Ausgaben auf der Kommandozeile.\\
Die Struktur der ledleuchtturm.py orientiert sich bis zum Aufruf der Funktion $F_9$ stark an dem Aktivitätsdiagramm. Ab dem ersten Ansteuern der LEDs verfolgt die Datei eine leicht veränderte Struktur.\\
Die Datei besitzt selbst nur die Funktion \texttt{signal\_handler}. Diese  wurde nicht ausgelagert, da sie auf globale Variablen aus mehreren Modulen zugreift. Deshalb ist eine zentrale Definition der Funktion von Vorteil.

\subsubsection*{fileFunc.py}
Das Modul fileFunc.py besitzt alle Funktionen die sich mit dem Erstellen, beschreiben oder lesen von Dateien beschäftigen.

\subsubsection*{dbFunc.py}
Das Modul dbFunc.py besitzt alle Funktionen die sich mit der Datenbank beschäftigen.

\subsubsection*{ledFunc.py}
Das Modul ledFunc.py besitzt alle Funktionen die sich mit der Ansteuerung oder Initialisierung der LEDs beschäftigen.

\subsubsection*{stdFunc.py}
Das Modul stdFunc.py besitzt Hilfsfunktionen die für den grundsätzlichen Programmablauf nicht wichtig sind. Die Funktionen darin werden im Entwicklungsprozess oder für die Fehlersuche verwendet.

% Muss so eindeutig sein, dass Fremde mithilfe dieses Abschnitts einen Überblick erhalten können und zusammen mit dem PyDoc das Programm ohne viel Nachforschung weiterentwickeln könnten => Anleitung, wie etwas verändert werden könnte!

% Struktur des Codes (wie er abgegeben wurde) festhalten. 
% Wie sind die Dateien/Packet(e) organisiert?
% Alle Dateien beschreiben: Überblick verschaffen.
% Am besten (statt Klassendiagramm) ein "Datei-Diagramm" erstellen.
% Muss so eindeutig sein, dass Fremde mithilfe dises Abschnitts einen Überblick erhalten können und zusammen mit dem PyDoc das Programm ohne viel Nachforschung weiterentwickeln könnten => Anleitung, wie etwas verändert werden könnte!

\clearpage
\section{EpyDoc}
\label{epydoc}
\begingroup
\setlength{\headheight}{16pt}
\setlength{\headsep}{24pt}
\setlength{\topmargin}{-\headsep}
\setlength{\parindent}{0ex}
\setlength{\parskip}{2ex}
\setlength{\fboxrule}{2\fboxrule}
\newlength{\BCL} % base class length, for base trees.
\newlength{\funcindent}
\newlength{\funcwidth}
\setlength{\funcindent}{1cm}
\setlength{\funcwidth}{\textwidth}
\addtolength{\funcwidth}{-2\funcindent}
\newlength{\varindent}
\newlength{\varnamewidth}
\newlength{\vardescrwidth}
\setlength{\varindent}{1cm}
\setlength{\varnamewidth}{.3\textwidth}
%\setlength{\varwidth}{\textwidth}
%\addtolength{\varwidth}{-4\tabcolsep}
%\addtolength{\varwidth}{-3\arrayrulewidth}
%\addtolength{\varwidth}{-2\varindent}
\setlength{\vardescrwidth}{8.3cm}
%\addtolength{\vardescrwidth}{-\varnamewidth}
\newenvironment{Ventry}[1]%
 {\begin{list}{}{%
   \renewcommand{\makelabel}[1]{\texttt{##1:}\hfil}%
   \settowidth{\labelwidth}{\texttt{#1:}}%
   \setlength{\leftmargin}{\labelsep}%
   \addtolength{\leftmargin}{\labelwidth}}}%
 {\end{list}}
\arrayrulecolor{black}
\LetLtxMacro{\oldsec}{\section}
\LetLtxMacro{\oldsubsec}{\subsection}
\renewcommand{\section}[1]{\oldsubsec{#1}}
\renewcommand{\subsection}[1]{\subsubsection{#1}}
%
% API Documentation for API Documentation
% Package ledleuchtturm
%
% Generated by epydoc 3.0.1
% [Tue Jul 18 00:55:24 2017]
%

%%%%%%%%%%%%%%%%%%%%%%%%%%%%%%%%%%%%%%%%%%%%%%%%%%%%%%%%%%%%%%%%%%%%%%%%%%%
%%                          Module Description                           %%
%%%%%%%%%%%%%%%%%%%%%%%%%%%%%%%%%%%%%%%%%%%%%%%%%%%%%%%%%%%%%%%%%%%%%%%%%%%

    \index{ledleuchtturm \textit{(package)}|(}
\section{Paket ledleuchtturm}

    \label{ledleuchtturm}
Das Paket ledleuchtturm verbindet sich mit einer postgresql Datenbank und 
wartet auf Aktualisierungen einer definierten Maschine. Bei einer 
Aktualisierung werden definierte Ports eines Raspberry Pi ein- oder 
ausgeschaltet. Direkt nach dem Verbindungsaufbau wird zuerst der aktuelle 
Status gelesen und dann werden die entsprechenden Pins geschaltet. Danach 
wird auf Updates gewartet. Durch das Signal Strg-C werden die Ports des 
Raspberry Pi in ihren Anfangszustand versetzt und alle Verbindungen oder 
geöffneten Dateien geschlossen.


%%%%%%%%%%%%%%%%%%%%%%%%%%%%%%%%%%%%%%%%%%%%%%%%%%%%%%%%%%%%%%%%%%%%%%%%%%%
%%                                Modules                                %%
%%%%%%%%%%%%%%%%%%%%%%%%%%%%%%%%%%%%%%%%%%%%%%%%%%%%%%%%%%%%%%%%%%%%%%%%%%%

\subsection{Module}

\begin{itemize}
\setlength{\parskip}{0ex}
\item \textbf{dbFunc}: Funktionen zur Arbeit mit der Datenbank.



  \textit{(Unterabschnitt \ref{ledleuchtturm:dbFunc}, S.~\pageref{ledleuchtturm:dbFunc})}

\item \textbf{fileFunc}: Funktionen zur Arbeit mit Dateien.



  \textit{(Unterabschnitt \ref{ledleuchtturm:fileFunc}, S.~\pageref{ledleuchtturm:fileFunc})}

\item \textbf{ledFunc}: Funktionen zur Arbeit mit LEDs.



  \textit{(Unterabschnitt \ref{ledleuchtturm:ledFunc}, S.~\pageref{ledleuchtturm:ledFunc})}

\item \textbf{ledleuchtturm}: Main-Funktion des Paketes ledleuchtturm.



  \textit{(Unterabschnitt \ref{ledleuchtturm:ledleuchtturm}, S.~\pageref{ledleuchtturm:ledleuchtturm})}

\item \textbf{stdFunc}: Funktionen, die beim Entwicklungsprozess benutzt werden.



  \textit{(Unterabschnitt \ref{ledleuchtturm:stdFunc}, S.~\pageref{ledleuchtturm:stdFunc})}

\end{itemize}


%%%%%%%%%%%%%%%%%%%%%%%%%%%%%%%%%%%%%%%%%%%%%%%%%%%%%%%%%%%%%%%%%%%%%%%%%%%
%%                               Variables                               %%
%%%%%%%%%%%%%%%%%%%%%%%%%%%%%%%%%%%%%%%%%%%%%%%%%%%%%%%%%%%%%%%%%%%%%%%%%%%

  \subsection{Variablen}

    \vspace{-1cm}
\hspace{\varindent}\begin{longtable}{|p{\varnamewidth}|p{\vardescrwidth}|l}
\cline{1-2}
\cline{1-2} \centering \textbf{Name} & \centering \textbf{Beschreibung}& \\
\cline{1-2}
\endhead\cline{1-2}\multicolumn{3}{r}{\small\ldots}\\\endfoot\cline{1-2}
\endlastfoot\raggedright \_\-\_\-p\-a\-c\-k\-a\-g\-e\-\_\-\_\- & \raggedright \textbf{Wert:} 
{\tt None}&\\
\cline{1-2}
\end{longtable}

    \index{ledleuchtturm \textit{(package)}|)}

%
% API Documentation for API Documentation
% Module ledleuchtturm.dbFunc
%
% Generated by epydoc 3.0.1
% [Tue Jul 18 00:55:24 2017]
%

%%%%%%%%%%%%%%%%%%%%%%%%%%%%%%%%%%%%%%%%%%%%%%%%%%%%%%%%%%%%%%%%%%%%%%%%%%%
%%                          Module Description                           %%
%%%%%%%%%%%%%%%%%%%%%%%%%%%%%%%%%%%%%%%%%%%%%%%%%%%%%%%%%%%%%%%%%%%%%%%%%%%

    \index{ledleuchtturm \textit{(package)}!ledleuchtturm.dbFunc \textit{(module)}|(}
\section{Modul ledleuchtturm.dbFunc}

    \label{ledleuchtturm:dbFunc}
Funktionen zur Arbeit mit der Datenbank.

\textbf{Autor:} Simon Retsch




%%%%%%%%%%%%%%%%%%%%%%%%%%%%%%%%%%%%%%%%%%%%%%%%%%%%%%%%%%%%%%%%%%%%%%%%%%%
%%                               Functions                               %%
%%%%%%%%%%%%%%%%%%%%%%%%%%%%%%%%%%%%%%%%%%%%%%%%%%%%%%%%%%%%%%%%%%%%%%%%%%%

  \subsection{Funktionen}

    \label{ledleuchtturm:dbFunc:testDBconnection}
    \index{ledleuchtturm \textit{(package)}!ledleuchtturm.dbFunc \textit{(module)}!ledleuchtturm.dbFunc.testDBconnection \textit{(function)}}

    \vspace{0.5ex}

\hspace{.8\funcindent}\begin{boxedminipage}{\funcwidth}

    \raggedright \textbf{testDBconnection}(\textit{configData})

    \vspace{-1.5ex}

    \rule{\textwidth}{0.5\fboxrule}
\setlength{\parskip}{2ex}
    Testen der Datenbankverbindung.

\setlength{\parskip}{1ex}
      \textbf{Parameter}
      \vspace{-1ex}

      \begin{quote}
        \begin{Ventry}{xxxxxxxxxx}

          \item[configData]

          Beinhaltet die Daten der Konfigurationsdatei.

            {\it (type=Json Object)}

        \end{Ventry}

      \end{quote}

      \textbf{R"uckgabewert}
    \vspace{-1ex}

      \begin{quote}
      Verbindungsvariable falls die Verbindung hergestellt werden kann, -1 
      falls die Verbindung nicht hergestellt werden kann.

      {\it (type=int / psycopg2 Connectionvariable)}

      \end{quote}

\textbf{Autor:} Simon Retsch



    \end{boxedminipage}

    \label{ledleuchtturm:dbFunc:waitForDBupdate}
    \index{ledleuchtturm \textit{(package)}!ledleuchtturm.dbFunc \textit{(module)}!ledleuchtturm.dbFunc.waitForDBupdate \textit{(function)}}

    \vspace{0.5ex}

\hspace{.8\funcindent}\begin{boxedminipage}{\funcwidth}

    \raggedright \textbf{waitForDBupdate}(\textit{configData})

    \vspace{-1.5ex}

    \rule{\textwidth}{0.5\fboxrule}
\setlength{\parskip}{2ex}
    Verbindet sich über Sockets zu dem in der Konfigurationsdatei 
    festgelegten Postgresql Server und wartet auf ein Update der ebenfalls 
    dort festgelegten Maschine, der Status nach diesem Update wird 
    zurückgegeben.

\setlength{\parskip}{1ex}
      \textbf{Parameter}
      \vspace{-1ex}

      \begin{quote}
        \begin{Ventry}{xxxxxxxxxx}

          \item[configData]

          Beinhaltet die Daten aus der Konfigurationsdatei der 
          Konfigurationsdatei.

            {\it (type=Json Object)}

        \end{Ventry}

      \end{quote}

      \textbf{R"uckgabewert}
    \vspace{-1ex}

      \begin{quote}
      Status der Maschine machineName bei einem Update.

      {\it (type=string)}

      \end{quote}

\textbf{Autor:} Ragnar Luga (kleine Änderungen durch Simon Retsch)



    \end{boxedminipage}

    \label{ledleuchtturm:dbFunc:getCurrentDBstate}
    \index{ledleuchtturm \textit{(package)}!ledleuchtturm.dbFunc \textit{(module)}!ledleuchtturm.dbFunc.getCurrentDBstate \textit{(function)}}

    \vspace{0.5ex}

\hspace{.8\funcindent}\begin{boxedminipage}{\funcwidth}

    \raggedright \textbf{getCurrentDBstate}(\textit{dbConnection}, \textit{configData})

    \vspace{-1.5ex}

    \rule{\textwidth}{0.5\fboxrule}
\setlength{\parskip}{2ex}
    Erstellt einen temporären Cursor für die Datenbank, führt darüber ein 
    select auf die Tabelle tool aus, liest das Ergebnis ein, untersucht ob 
    ein valider Zustand für die in configData angegebene Maschine darin 
    existiert und gibt diesen zurück.

\setlength{\parskip}{1ex}
      \textbf{Parameter}
      \vspace{-1ex}

      \begin{quote}
        \begin{Ventry}{xxxxxxxxxxxx}

          \item[dbConnection]

          Variable der Verbindung mit der Datenbank.

            {\it (type=psycopg2 Connectionvariable)}

          \item[configData]

          Beinhaltet die Daten aus der Konfigurationsdatei der 
          Konfigurationsdatei.

            {\it (type=Json Object)}

        \end{Ventry}

      \end{quote}

      \textbf{R"uckgabewert}
    \vspace{-1ex}

      \begin{quote}
      Status der Maschine falls die Maschine gefunden wurde und ein valider
      Status vorhanden war, -1 falls nicht.

      {\it (type=int / string)}

      \end{quote}

\textbf{Autor:} Simon Retsch



    \end{boxedminipage}


%%%%%%%%%%%%%%%%%%%%%%%%%%%%%%%%%%%%%%%%%%%%%%%%%%%%%%%%%%%%%%%%%%%%%%%%%%%
%%                               Variables                               %%
%%%%%%%%%%%%%%%%%%%%%%%%%%%%%%%%%%%%%%%%%%%%%%%%%%%%%%%%%%%%%%%%%%%%%%%%%%%

  \subsection{Variablen}

    \vspace{-1cm}
\hspace{\varindent}\begin{longtable}{|p{\varnamewidth}|p{\vardescrwidth}|l}
\cline{1-2}
\cline{1-2} \centering \textbf{Name} & \centering \textbf{Beschreibung}& \\
\cline{1-2}
\endhead\cline{1-2}\multicolumn{3}{r}{\small\ldots}\\\endfoot\cline{1-2}
\endlastfoot\raggedright g\-l\-o\-b\-a\-l\-D\-B\-c\-o\-n\-n\-e\-c\-t\-i\-o\-n\- & \raggedright Verbindungsvariable zur Datenbank.

\textbf{Wert:} 
{\tt -1}            {\it (type=psycopg2 database connection object)}&\\
\cline{1-2}
\raggedright g\-l\-o\-b\-a\-l\-D\-B\-s\-o\-c\-k\-e\-t\- & \raggedright Socketvariable zur Datenbank.

\textbf{Wert:} 
{\tt -1}            {\it (type=psycopg2 socket object)}&\\
\cline{1-2}
\raggedright g\-l\-o\-b\-a\-l\-D\-B\-c\-u\-r\-s\-o\-r\- & \raggedright Cursorvariable zur Datenbank.

\textbf{Wert:} 
{\tt -1}            {\it (type=psycopg2 cursor object)}&\\
\cline{1-2}
\end{longtable}

    \index{ledleuchtturm \textit{(package)}!ledleuchtturm.dbFunc \textit{(module)}|)}

%
% API Documentation for API Documentation
% Module ledleuchtturm.fileFunc
%
% Generated by epydoc 3.0.1
% [Tue Jul 18 00:55:24 2017]
%

%%%%%%%%%%%%%%%%%%%%%%%%%%%%%%%%%%%%%%%%%%%%%%%%%%%%%%%%%%%%%%%%%%%%%%%%%%%
%%                          Module Description                           %%
%%%%%%%%%%%%%%%%%%%%%%%%%%%%%%%%%%%%%%%%%%%%%%%%%%%%%%%%%%%%%%%%%%%%%%%%%%%

    \index{ledleuchtturm \textit{(package)}!ledleuchtturm.fileFunc \textit{(module)}|(}
\section{Modul ledleuchtturm.fileFunc}

    \label{ledleuchtturm:fileFunc}
Funktionen zur Arbeit mit Dateien.

\textbf{Autor:} Simon Retsch




%%%%%%%%%%%%%%%%%%%%%%%%%%%%%%%%%%%%%%%%%%%%%%%%%%%%%%%%%%%%%%%%%%%%%%%%%%%
%%                               Functions                               %%
%%%%%%%%%%%%%%%%%%%%%%%%%%%%%%%%%%%%%%%%%%%%%%%%%%%%%%%%%%%%%%%%%%%%%%%%%%%

  \subsection{Funktionen}

    \label{ledleuchtturm:fileFunc:checkArgumentForConfigFile}
    \index{ledleuchtturm \textit{(package)}!ledleuchtturm.fileFunc \textit{(module)}!ledleuchtturm.fileFunc.checkArgumentForConfigFile \textit{(function)}}

    \vspace{0.5ex}

\hspace{.8\funcindent}\begin{boxedminipage}{\funcwidth}

    \raggedright \textbf{checkArgumentForConfigFile}(\textit{argv})

    \vspace{-1.5ex}

    \rule{\textwidth}{0.5\fboxrule}
\setlength{\parskip}{2ex}
    Überprüft ob dem Programm ein Pfad einer Konfigurationsdatei mit der 
    Endung .json als Argument übergeben wurde.

\setlength{\parskip}{1ex}
      \textbf{Parameter}
      \vspace{-1ex}

      \begin{quote}
        \begin{Ventry}{xxxx}

          \item[argv]

          Programm Übergabeparameter.

            {\it (type=string[])}

        \end{Ventry}

      \end{quote}

      \textbf{R"uckgabewert}
    \vspace{-1ex}

      \begin{quote}
      Pfad der Konfigurationsdatei falls es erfolgreich ist, -1 falls nicht
      die richtige Anzahl von Argumenten mitgegeben wurde.

      {\it (type=int / string)}

      \end{quote}

\textbf{Autor:} Raphael Pour



    \end{boxedminipage}

    \label{ledleuchtturm:fileFunc:configFileExists}
    \index{ledleuchtturm \textit{(package)}!ledleuchtturm.fileFunc \textit{(module)}!ledleuchtturm.fileFunc.configFileExists \textit{(function)}}

    \vspace{0.5ex}

\hspace{.8\funcindent}\begin{boxedminipage}{\funcwidth}

    \raggedright \textbf{configFileExists}(\textit{configFile})

    \vspace{-1.5ex}

    \rule{\textwidth}{0.5\fboxrule}
\setlength{\parskip}{2ex}
    Überprüft ob die Konfigurationsdatei existiert.

\setlength{\parskip}{1ex}
      \textbf{Parameter}
      \vspace{-1ex}

      \begin{quote}
        \begin{Ventry}{xxxxxxxxxx}

          \item[configFile]

          Pfad der Konfigurationsdatei.

            {\it (type=string)}

        \end{Ventry}

      \end{quote}

      \textbf{R"uckgabewert}
    \vspace{-1ex}

      \begin{quote}
      1 falls die Datei existiert, -1 falls die Datei nicht existiert.

      {\it (type=int)}

      \end{quote}

\textbf{Autor:} Raphael Pour



    \end{boxedminipage}

    \label{ledleuchtturm:fileFunc:configFileValidJsonFile}
    \index{ledleuchtturm \textit{(package)}!ledleuchtturm.fileFunc \textit{(module)}!ledleuchtturm.fileFunc.configFileValidJsonFile \textit{(function)}}

    \vspace{0.5ex}

\hspace{.8\funcindent}\begin{boxedminipage}{\funcwidth}

    \raggedright \textbf{configFileValidJsonFile}(\textit{configFile})

    \vspace{-1.5ex}

    \rule{\textwidth}{0.5\fboxrule}
\setlength{\parskip}{2ex}
    Überprüft ob es sich bei configFile wirklich um eine json Datei 
    handelt.

\setlength{\parskip}{1ex}
      \textbf{Parameter}
      \vspace{-1ex}

      \begin{quote}
        \begin{Ventry}{xxxxxxxxxx}

          \item[configFile]

          Pfad der Konfigurationsdatei.

            {\it (type=string)}

        \end{Ventry}

      \end{quote}

      \textbf{R"uckgabewert}
    \vspace{-1ex}

      \begin{quote}
      1 falls configFile eine Json Datei ist, -1 falls configFile keine 
      Json Datei ist.

      {\it (type=int)}

      \end{quote}

\textbf{Autor:} Raphael Pour



    \end{boxedminipage}

    \label{ledleuchtturm:fileFunc:createStdConf}
    \index{ledleuchtturm \textit{(package)}!ledleuchtturm.fileFunc \textit{(module)}!ledleuchtturm.fileFunc.createStdConf \textit{(function)}}

    \vspace{0.5ex}

\hspace{.8\funcindent}\begin{boxedminipage}{\funcwidth}

    \raggedright \textbf{createStdConf}()

    \vspace{-1.5ex}

    \rule{\textwidth}{0.5\fboxrule}
\setlength{\parskip}{2ex}
    Schreibt ein Python Objekt in eine Datei mit dem Namen 
    Std-Conf-Ledleuchtturm.json. In dem Objekt befinden sich Standartwerte.

\setlength{\parskip}{1ex}
      \textbf{R"uckgabewert}
    \vspace{-1ex}

      \begin{quote}
      1, falls das Objekt in die Datei geschrieben werden kann, -1 falls es
      dabei Probleme gibt.

      {\it (type=int)}

      \end{quote}

\textbf{Autor:} Simon Retsch



    \end{boxedminipage}

    \label{ledleuchtturm:fileFunc:getJsonFileObject}
    \index{ledleuchtturm \textit{(package)}!ledleuchtturm.fileFunc \textit{(module)}!ledleuchtturm.fileFunc.getJsonFileObject \textit{(function)}}

    \vspace{0.5ex}

\hspace{.8\funcindent}\begin{boxedminipage}{\funcwidth}

    \raggedright \textbf{getJsonFileObject}(\textit{fileName})

    \vspace{-1.5ex}

    \rule{\textwidth}{0.5\fboxrule}
\setlength{\parskip}{2ex}
    Lädt die Struktur einer Json Datei in ein Objekt.

\setlength{\parskip}{1ex}
      \textbf{Parameter}
      \vspace{-1ex}

      \begin{quote}
        \begin{Ventry}{xxxxxxxx}

          \item[fileName]

          Pfad der Json Datei.

            {\it (type=string)}

        \end{Ventry}

      \end{quote}

      \textbf{R"uckgabewert}
    \vspace{-1ex}

      \begin{quote}
      Objekt welches aus einer Json Datei geladen wurde.

      {\it (type=Json Object)}

      \end{quote}

\textbf{Autor:} Simon Retsch



    \end{boxedminipage}


%%%%%%%%%%%%%%%%%%%%%%%%%%%%%%%%%%%%%%%%%%%%%%%%%%%%%%%%%%%%%%%%%%%%%%%%%%%
%%                               Variables                               %%
%%%%%%%%%%%%%%%%%%%%%%%%%%%%%%%%%%%%%%%%%%%%%%%%%%%%%%%%%%%%%%%%%%%%%%%%%%%

  \subsection{Variablen}

    \vspace{-1cm}
\hspace{\varindent}\begin{longtable}{|p{\varnamewidth}|p{\vardescrwidth}|l}
\cline{1-2}
\cline{1-2} \centering \textbf{Name} & \centering \textbf{Beschreibung}& \\
\cline{1-2}
\endhead\cline{1-2}\multicolumn{3}{r}{\small\ldots}\\\endfoot\cline{1-2}
\endlastfoot\raggedright g\-l\-o\-b\-a\-l\-F\-i\-l\-e\-O\-p\-e\-n\- & \raggedright Dateiobjekt einer geöffneten Datei.

\textbf{Wert:} 
{\tt -1}            {\it (type=File object)}&\\
\cline{1-2}
\raggedright \_\-\_\-p\-a\-c\-k\-a\-g\-e\-\_\-\_\- & \raggedright \textbf{Wert:} 
{\tt \texttt{'}\texttt{ledleuchtturm}\texttt{'}}&\\
\cline{1-2}
\end{longtable}

    \index{ledleuchtturm \textit{(package)}!ledleuchtturm.fileFunc \textit{(module)}|)}

%
% API Documentation for API Documentation
% Module ledleuchtturm.ledFunc
%
% Generated by epydoc 3.0.1
% [Tue Jul 18 00:55:24 2017]
%

%%%%%%%%%%%%%%%%%%%%%%%%%%%%%%%%%%%%%%%%%%%%%%%%%%%%%%%%%%%%%%%%%%%%%%%%%%%
%%                          Module Description                           %%
%%%%%%%%%%%%%%%%%%%%%%%%%%%%%%%%%%%%%%%%%%%%%%%%%%%%%%%%%%%%%%%%%%%%%%%%%%%

    \index{ledleuchtturm \textit{(package)}!ledleuchtturm.ledFunc \textit{(module)}|(}
\section{Modul ledleuchtturm.ledFunc}

    \label{ledleuchtturm:ledFunc}
Funktionen zur Arbeit mit LEDs.

\textbf{Autor:} Simon Retsch




%%%%%%%%%%%%%%%%%%%%%%%%%%%%%%%%%%%%%%%%%%%%%%%%%%%%%%%%%%%%%%%%%%%%%%%%%%%
%%                               Functions                               %%
%%%%%%%%%%%%%%%%%%%%%%%%%%%%%%%%%%%%%%%%%%%%%%%%%%%%%%%%%%%%%%%%%%%%%%%%%%%

  \subsection{Funktionen}

    \label{ledleuchtturm:ledFunc:toggleLED}
    \index{ledleuchtturm \textit{(package)}!ledleuchtturm.ledFunc \textit{(module)}!ledleuchtturm.ledFunc.toggleLED \textit{(function)}}

    \vspace{0.5ex}

\hspace{.8\funcindent}\begin{boxedminipage}{\funcwidth}

    \raggedright \textbf{toggleLED}(\textit{portnumber})

    \vspace{-1.5ex}

    \rule{\textwidth}{0.5\fboxrule}
\setlength{\parskip}{2ex}
    Wechselt die den Zustand (an/aus) der LED am angegebenen Port.

\setlength{\parskip}{1ex}
      \textbf{Parameter}
      \vspace{-1ex}

      \begin{quote}
        \begin{Ventry}{xxxxxxxxxx}

          \item[portnumber]

          Nummer des Ports am Rasperry Pi.

            {\it (type=int)}

        \end{Ventry}

      \end{quote}

      \textbf{R"uckgabewert}
    \vspace{-1ex}

      \begin{quote}
      1 bei Erfolg, -1 bei Misserfolg.

      {\it (type=int)}

      \end{quote}

\textbf{Autor:} Falk-Jonatan Strube



    \end{boxedminipage}

    \label{ledleuchtturm:ledFunc:enableLED}
    \index{ledleuchtturm \textit{(package)}!ledleuchtturm.ledFunc \textit{(module)}!ledleuchtturm.ledFunc.enableLED \textit{(function)}}

    \vspace{0.5ex}

\hspace{.8\funcindent}\begin{boxedminipage}{\funcwidth}

    \raggedright \textbf{enableLED}(\textit{portnumber})

    \vspace{-1.5ex}

    \rule{\textwidth}{0.5\fboxrule}
\setlength{\parskip}{2ex}
    Schaltet die LED am angegeben Port an.

\setlength{\parskip}{1ex}
      \textbf{Parameter}
      \vspace{-1ex}

      \begin{quote}
        \begin{Ventry}{xxxxxxxxxx}

          \item[portnumber]

          Nummer des Ports am Rasperry Pi.

            {\it (type=int)}

        \end{Ventry}

      \end{quote}

      \textbf{R"uckgabewert}
    \vspace{-1ex}

      \begin{quote}
      1 bei Erfolg, -1 bei Misserfolg.

      {\it (type=int)}

      \end{quote}

\textbf{Autor:} Falk-Jonatan Strube



    \end{boxedminipage}

    \label{ledleuchtturm:ledFunc:switchAllOff}
    \index{ledleuchtturm \textit{(package)}!ledleuchtturm.ledFunc \textit{(module)}!ledleuchtturm.ledFunc.switchAllOff \textit{(function)}}

    \vspace{0.5ex}

\hspace{.8\funcindent}\begin{boxedminipage}{\funcwidth}

    \raggedright \textbf{switchAllOff}()

    \vspace{-1.5ex}

    \rule{\textwidth}{0.5\fboxrule}
\setlength{\parskip}{2ex}
    Schaltet alle LEDs (definiert durch die globale Variablen der 
    entsprechenden Ports) aus.

\setlength{\parskip}{1ex}
      \textbf{R"uckgabewert}
    \vspace{-1ex}

      \begin{quote}
      1 bei Erfolg, -1 bei Misserfolg.

      {\it (type=int)}

      \end{quote}

\textbf{Autor:} Falk-Jonatan Strube



    \end{boxedminipage}

    \label{ledleuchtturm:ledFunc:blink}
    \index{ledleuchtturm \textit{(package)}!ledleuchtturm.ledFunc \textit{(module)}!ledleuchtturm.ledFunc.blink \textit{(function)}}

    \vspace{0.5ex}

\hspace{.8\funcindent}\begin{boxedminipage}{\funcwidth}

    \raggedright \textbf{blink}(\textit{portnumber})

    \vspace{-1.5ex}

    \rule{\textwidth}{0.5\fboxrule}
\setlength{\parskip}{2ex}
    Lässt die LED am angegebenen Port blinken, indem der ein Thread in der 
    globalen Variabel \texttt{blinkThread} als \texttt{blinkThreadFunction}
    gestartet wird.

\setlength{\parskip}{1ex}
      \textbf{Parameter}
      \vspace{-1ex}

      \begin{quote}
        \begin{Ventry}{xxxxxxxxxx}

          \item[portnumber]

          Nummer des Ports am Rasperry Pi.

            {\it (type=int)}

        \end{Ventry}

      \end{quote}

      \textbf{R"uckgabewert}
    \vspace{-1ex}

      \begin{quote}
      1 bei Erfolg, -1 bei Misserfolg.

      {\it (type=int)}

      \end{quote}

\textbf{Autor:} Falk-Jonatan Strube



    \end{boxedminipage}

    \label{ledleuchtturm:ledFunc:blinkThreadFunction}
    \index{ledleuchtturm \textit{(package)}!ledleuchtturm.ledFunc \textit{(module)}!ledleuchtturm.ledFunc.blinkThreadFunction \textit{(function)}}

    \vspace{0.5ex}

\hspace{.8\funcindent}\begin{boxedminipage}{\funcwidth}

    \raggedright \textbf{blinkThreadFunction}(\textit{portnumber})

    \vspace{-1.5ex}

    \rule{\textwidth}{0.5\fboxrule}
\setlength{\parskip}{2ex}
    Diese Funktion wird als Thread genutzt und lässt die LED am angegeben 
    Port blinken: Das Warten zwischen dem Wechseln des LED-Zustands wird 
    über das globale \texttt{blinkThreadStopEvent} realisiert. Dadurch kann
    der Zustand des Blinkens jederzeit beendet werden.

\setlength{\parskip}{1ex}
      \textbf{Parameter}
      \vspace{-1ex}

      \begin{quote}
        \begin{Ventry}{xxxxxxxxxx}

          \item[portnumber]

          Nummer des Ports am Rasperry Pi.

            {\it (type=int)}

        \end{Ventry}

      \end{quote}

      \textbf{R"uckgabewert}
    \vspace{-1ex}

      \begin{quote}
      1 bei Erfolg, -1 bei Misserfolg.

      {\it (type=int)}

      \end{quote}

\textbf{Autor:} Falk-Jonatan Strube



    \end{boxedminipage}

    \label{ledleuchtturm:ledFunc:killBlinkThread}
    \index{ledleuchtturm \textit{(package)}!ledleuchtturm.ledFunc \textit{(module)}!ledleuchtturm.ledFunc.killBlinkThread \textit{(function)}}

    \vspace{0.5ex}

\hspace{.8\funcindent}\begin{boxedminipage}{\funcwidth}

    \raggedright \textbf{killBlinkThread}()

    \vspace{-1.5ex}

    \rule{\textwidth}{0.5\fboxrule}
\setlength{\parskip}{2ex}
    Beendet den für das Blinken verantwortliche Thread 
    \texttt{blinkThread}. Es wird zuerst das \texttt{blinkThreadStopEvent} 
    aktiviert und der Thread dann zusammengeführt.

\setlength{\parskip}{1ex}
      \textbf{R"uckgabewert}
    \vspace{-1ex}

      \begin{quote}
      1 bei Erfolg, -1 bei Misserfolg.

      {\it (type=int)}

      \end{quote}

\textbf{Autor:} Falk-Jonatan Strube



    \end{boxedminipage}

    \label{ledleuchtturm:ledFunc:changeState}
    \index{ledleuchtturm \textit{(package)}!ledleuchtturm.ledFunc \textit{(module)}!ledleuchtturm.ledFunc.changeState \textit{(function)}}

    \vspace{0.5ex}

\hspace{.8\funcindent}\begin{boxedminipage}{\funcwidth}

    \raggedright \textbf{changeState}(\textit{status})

    \vspace{-1.5ex}

    \rule{\textwidth}{0.5\fboxrule}
\setlength{\parskip}{2ex}
    Wechselt den Status der LED-Anzeige. Dabei werden zuerst Blink-Threads 
    beendet (falls vorhanden). Dann wird die dem übergebenen Status 
    entsprechende LED eingeschaltet oder das Blinken gestartet. Wenn der 
    übergebene Status unbekannt ist, werden alle LEDs ausgeschaltet.

\setlength{\parskip}{1ex}
      \textbf{Parameter}
      \vspace{-1ex}

      \begin{quote}
        \begin{Ventry}{xxxxxx}

          \item[status]

          Status, zu dem gewechselt werden soll (bekannte Stati: 
          PROC/IDLE/MAINT/DOWN).

            {\it (type=string)}

        \end{Ventry}

      \end{quote}

      \textbf{R"uckgabewert}
    \vspace{-1ex}

      \begin{quote}
      1 bei Erfolg, -1 bei Misserfolg.

      {\it (type=int)}

      \end{quote}

\textbf{Autor:} Falk-Jonatan Strube



    \end{boxedminipage}

    \label{ledleuchtturm:ledFunc:setLEDsGlobal}
    \index{ledleuchtturm \textit{(package)}!ledleuchtturm.ledFunc \textit{(module)}!ledleuchtturm.ledFunc.setLEDsGlobal \textit{(function)}}

    \vspace{0.5ex}

\hspace{.8\funcindent}\begin{boxedminipage}{\funcwidth}

    \raggedright \textbf{setLEDsGlobal}(\textit{configData})

    \vspace{-1.5ex}

    \rule{\textwidth}{0.5\fboxrule}
\setlength{\parskip}{2ex}
    Setzt die Globalen Variablen für die LED Pins mit Hilfe der Daten aus 
    der Konfigurationsdatei.

\setlength{\parskip}{1ex}
      \textbf{Parameter}
      \vspace{-1ex}

      \begin{quote}
        \begin{Ventry}{xxxxxxxxxx}

          \item[configData]

          Beinhaltet die Daten aus der Konfigurationsdatei der 
          Konfigurationsdatei.

            {\it (type=Json Object)}

        \end{Ventry}

      \end{quote}

\textbf{Autor:} Simon Retsch



    \end{boxedminipage}

    \label{ledleuchtturm:ledFunc:initLEDs}
    \index{ledleuchtturm \textit{(package)}!ledleuchtturm.ledFunc \textit{(module)}!ledleuchtturm.ledFunc.initLEDs \textit{(function)}}

    \vspace{0.5ex}

\hspace{.8\funcindent}\begin{boxedminipage}{\funcwidth}

    \raggedright \textbf{initLEDs}()

    \vspace{-1.5ex}

    \rule{\textwidth}{0.5\fboxrule}
\setlength{\parskip}{2ex}
    Initialisiert die LEDs des Raspberry Pi: Setzen des Pinning-Layouts und
    Pins als Output definieren

\setlength{\parskip}{1ex}
      \textbf{R"uckgabewert}
    \vspace{-1ex}

      \begin{quote}
      1 falls das Initialisieren erfolgreich war, -1 wenn nicht.

      {\it (type=int)}

      \end{quote}

\textbf{Autor:} Simon Retsch



    \end{boxedminipage}


%%%%%%%%%%%%%%%%%%%%%%%%%%%%%%%%%%%%%%%%%%%%%%%%%%%%%%%%%%%%%%%%%%%%%%%%%%%
%%                               Variables                               %%
%%%%%%%%%%%%%%%%%%%%%%%%%%%%%%%%%%%%%%%%%%%%%%%%%%%%%%%%%%%%%%%%%%%%%%%%%%%

  \subsection{Variablen}

    \vspace{-1cm}
\hspace{\varindent}\begin{longtable}{|p{\varnamewidth}|p{\vardescrwidth}|l}
\cline{1-2}
\cline{1-2} \centering \textbf{Name} & \centering \textbf{Beschreibung}& \\
\cline{1-2}
\endhead\cline{1-2}\multicolumn{3}{r}{\small\ldots}\\\endfoot\cline{1-2}
\endlastfoot\raggedright G\-R\-E\-E\-N\-\_\-P\-O\-R\-T\- & \raggedright Port der grünen LED auf dem Raspberry Pi.

\textbf{Wert:} 
{\tt -1}            {\it (type=int)}&\\
\cline{1-2}
\raggedright Y\-E\-L\-L\-O\-W\-\_\-P\-O\-R\-T\- & \raggedright Port der gelben LED auf dem Raspberry Pi.

\textbf{Wert:} 
{\tt -1}            {\it (type=int)}&\\
\cline{1-2}
\raggedright R\-E\-D\-\_\-P\-O\-R\-T\- & \raggedright Port der roten LED auf dem Raspberry Pi.

\textbf{Wert:} 
{\tt -1}            {\it (type=int)}&\\
\cline{1-2}
\raggedright b\-l\-i\-n\-k\-T\-h\-r\-e\-a\-d\- & \raggedright Objekt des Threads der das Blinken einer LED realisiert.

\textbf{Wert:} 
{\tt threading.Thread()}            {\it (type=Thread object)}&\\
\cline{1-2}
\raggedright b\-l\-i\-n\-k\-T\-h\-r\-e\-a\-d\-S\-t\-o\-p\-E\-v\-e\-n\-t\- & \raggedright Objekt eines Thread Events, welches den Thread beendet.

\textbf{Wert:} 
{\tt threading.Event()}            {\it (type=Thread event object)}&\\
\cline{1-2}
\end{longtable}

    \index{ledleuchtturm \textit{(package)}!ledleuchtturm.ledFunc \textit{(module)}|)}

%
% API Documentation for API Documentation
% Module ledleuchtturm.ledleuchtturm
%
% Generated by epydoc 3.0.1
% [Tue Jul 18 00:55:24 2017]
%

%%%%%%%%%%%%%%%%%%%%%%%%%%%%%%%%%%%%%%%%%%%%%%%%%%%%%%%%%%%%%%%%%%%%%%%%%%%
%%                          Module Description                           %%
%%%%%%%%%%%%%%%%%%%%%%%%%%%%%%%%%%%%%%%%%%%%%%%%%%%%%%%%%%%%%%%%%%%%%%%%%%%

    \index{ledleuchtturm \textit{(package)}!ledleuchtturm.ledleuchtturm \textit{(module)}|(}
\section{Modul ledleuchtturm.ledleuchtturm}

    \label{ledleuchtturm:ledleuchtturm}
Main-Funktion des Paketes ledleuchtturm. Sie repräsentiert den Ablauf des 
Programms. Sie beschränkt sich auf das Aufrufen von Funktionen aus den 
Modulen \texttt{stdFunc}, \texttt{fileFunc}, \texttt{dbFunc} oder 
\texttt{ledFunc}. Des Weiteren behandelt  sie die Auswertung von 
Rückgabewerten und das Schreiben von Meldungen auf der Konsole.

\textbf{Autor:} Simon Retsch




%%%%%%%%%%%%%%%%%%%%%%%%%%%%%%%%%%%%%%%%%%%%%%%%%%%%%%%%%%%%%%%%%%%%%%%%%%%
%%                               Functions                               %%
%%%%%%%%%%%%%%%%%%%%%%%%%%%%%%%%%%%%%%%%%%%%%%%%%%%%%%%%%%%%%%%%%%%%%%%%%%%

  \subsection{Funktionen}

    \label{ledleuchtturm:ledleuchtturm:signal_handler}
    \index{ledleuchtturm \textit{(package)}!ledleuchtturm.ledleuchtturm \textit{(module)}!ledleuchtturm.ledleuchtturm.signal\_handler \textit{(function)}}

    \vspace{0.5ex}

\hspace{.8\funcindent}\begin{boxedminipage}{\funcwidth}

    \raggedright \textbf{signal\_handler}(\textit{signal}, \textit{frame})

    \vspace{-1.5ex}

    \rule{\textwidth}{0.5\fboxrule}
\setlength{\parskip}{2ex}
    Behandelt das Signal welches durch Strg+C ausgelöst wird (SIGINT) und 
    beendet alle Verbindungen oder schließt alle offenen Dateien damit das 
    Programm ordnungsgemäß beendet wird.

\setlength{\parskip}{1ex}
      \textbf{Parameter}
      \vspace{-1ex}

      \begin{quote}
        \begin{Ventry}{xxxxxx}

          \item[signal]

          Signalnummer

            {\it (type=int)}

          \item[frame]

          Unterbrochener Stack Frame

            {\it (type=stack frame)}

        \end{Ventry}

      \end{quote}

\textbf{Autor:} Simon Retsch



    \end{boxedminipage}


%%%%%%%%%%%%%%%%%%%%%%%%%%%%%%%%%%%%%%%%%%%%%%%%%%%%%%%%%%%%%%%%%%%%%%%%%%%
%%                               Variables                               %%
%%%%%%%%%%%%%%%%%%%%%%%%%%%%%%%%%%%%%%%%%%%%%%%%%%%%%%%%%%%%%%%%%%%%%%%%%%%

  \subsection{Variablen}

    \vspace{-1cm}
\hspace{\varindent}\begin{longtable}{|p{\varnamewidth}|p{\vardescrwidth}|l}
\cline{1-2}
\cline{1-2} \centering \textbf{Name} & \centering \textbf{Beschreibung}& \\
\cline{1-2}
\endhead\cline{1-2}\multicolumn{3}{r}{\small\ldots}\\\endfoot\cline{1-2}
\endlastfoot\raggedright D\-E\-B\-U\-G\- & \raggedright Entscheidung über Ausgabe der debug-Nachrichten.

\textbf{Wert:} 
{\tt 1}            {\it (type=int)}&\\
\cline{1-2}
\end{longtable}

    \index{ledleuchtturm \textit{(package)}!ledleuchtturm.ledleuchtturm \textit{(module)}|)}

%
% API Documentation for API Documentation
% Module ledleuchtturm.stdFunc
%
% Generated by epydoc 3.0.1
% [Tue Jul 18 00:55:24 2017]
%

%%%%%%%%%%%%%%%%%%%%%%%%%%%%%%%%%%%%%%%%%%%%%%%%%%%%%%%%%%%%%%%%%%%%%%%%%%%
%%                          Module Description                           %%
%%%%%%%%%%%%%%%%%%%%%%%%%%%%%%%%%%%%%%%%%%%%%%%%%%%%%%%%%%%%%%%%%%%%%%%%%%%

    \index{ledleuchtturm \textit{(package)}!ledleuchtturm.stdFunc \textit{(module)}|(}
\section{Modul ledleuchtturm.stdFunc}

    \label{ledleuchtturm:stdFunc}
Funktionen, die beim Entwicklungsprozess benutzt werden.

\textbf{Autor:} Simon Retsch




%%%%%%%%%%%%%%%%%%%%%%%%%%%%%%%%%%%%%%%%%%%%%%%%%%%%%%%%%%%%%%%%%%%%%%%%%%%
%%                               Functions                               %%
%%%%%%%%%%%%%%%%%%%%%%%%%%%%%%%%%%%%%%%%%%%%%%%%%%%%%%%%%%%%%%%%%%%%%%%%%%%

  \subsection{Funktionen}

    \label{ledleuchtturm:stdFunc:printConfigSettings}
    \index{ledleuchtturm \textit{(package)}!ledleuchtturm.stdFunc \textit{(module)}!ledleuchtturm.stdFunc.printConfigSettings \textit{(function)}}

    \vspace{0.5ex}

\hspace{.8\funcindent}\begin{boxedminipage}{\funcwidth}

    \raggedright \textbf{printConfigSettings}(\textit{jsonFileObject})

    \vspace{-1.5ex}

    \rule{\textwidth}{0.5\fboxrule}
\setlength{\parskip}{2ex}
    Gibt den Inhalt des jsonFileObjects auf der Konsole aus.

\setlength{\parskip}{1ex}
      \textbf{Parameter}
      \vspace{-1ex}

      \begin{quote}
        \begin{Ventry}{xxxxxxxxxxxxxx}

          \item[jsonFileObject]

          Objekt welches Daten der Konfigurationsdatei enthält.

            {\it (type=Json Object)}

        \end{Ventry}

      \end{quote}

\textbf{Autor:} Simon Retsch



    \end{boxedminipage}

    \label{ledleuchtturm:stdFunc:debug}
    \index{ledleuchtturm \textit{(package)}!ledleuchtturm.stdFunc \textit{(module)}!ledleuchtturm.stdFunc.debug \textit{(function)}}

    \vspace{0.5ex}

\hspace{.8\funcindent}\begin{boxedminipage}{\funcwidth}

    \raggedright \textbf{debug}(\textit{string}, \textit{DEBUG})

    \vspace{-1.5ex}

    \rule{\textwidth}{0.5\fboxrule}
\setlength{\parskip}{2ex}
    Schreibt die Nachricht string auf der Konsole, falls DEBUG auf 1 
    gesetzt ist.

\setlength{\parskip}{1ex}
      \textbf{Parameter}
      \vspace{-1ex}

      \begin{quote}
        \begin{Ventry}{xxxxxx}

          \item[string]

          Nachricht die auf der Konsole erscheinen soll.

            {\it (type=string)}

          \item[DEBUG]

          Abhängigkeitsvariable.

            {\it (type=int)}

        \end{Ventry}

      \end{quote}

\textbf{Autor:} Simon Retsch



    \end{boxedminipage}


%%%%%%%%%%%%%%%%%%%%%%%%%%%%%%%%%%%%%%%%%%%%%%%%%%%%%%%%%%%%%%%%%%%%%%%%%%%
%%                               Variables                               %%
%%%%%%%%%%%%%%%%%%%%%%%%%%%%%%%%%%%%%%%%%%%%%%%%%%%%%%%%%%%%%%%%%%%%%%%%%%%

  \subsection{Variablen}

    \vspace{-1cm}
\hspace{\varindent}\begin{longtable}{|p{\varnamewidth}|p{\vardescrwidth}|l}
\cline{1-2}
\cline{1-2} \centering \textbf{Name} & \centering \textbf{Beschreibung}& \\
\cline{1-2}
\endhead\cline{1-2}\multicolumn{3}{r}{\small\ldots}\\\endfoot\cline{1-2}
\endlastfoot\raggedright \_\-\_\-p\-a\-c\-k\-a\-g\-e\-\_\-\_\- & \raggedright \textbf{Wert:} 
{\tt None}&\\
\cline{1-2}
\end{longtable}

    \index{ledleuchtturm \textit{(package)}!ledleuchtturm.stdFunc \textit{(module)}|)}

\endgroup
% wird zum Schluss generiert und eingebunden

\end{document}